\chapter{Overview of the proofs}\label{sec:overview}

This section gives an overview of the proofs of \cref{thm:barrier-730} and
\cref{thm:upper}, and a shorter account of the parts of the paper that concern
\cref{f:main} and \cref{cor:mean-bias}. Every claim made here is proved later in
the paper, or quoted there with its source.
\Cref{sec:ov-sign,sec:ov-architecture,sec:ov-scales} explain the reduction to
paths, the architecture of the proof in \cite{oai-barriers}, which we follow,
and the balance of scales that gives the exponent $7/30$.
\Cref{sec:ov-cubic,sec:ov-static,sec:ov-dynamics} describe the cubic pressure
deficit (\cref{sec:scalar,sec:flow,sec:cubic}), the passage to Gibbs
probabilities (\cref{sec:fractional,sec:static}), and the dynamical argument
(\cref{sec:references,sec:induction,sec:barrier-proof}). \Cref{sec:ov-upper}
concerns \cref{part:upper}, and \cref{sec:ov-fluct} concerns
\cref{part:fluct}. \Cref{tab:ov-dependencies} shows how the sections depend on
each other.

Throughout this section, $\beta>1$ is fixed. In
\cref{sec:ov-sign,sec:ov-architecture,sec:ov-scales,sec:ov-cubic,sec:ov-static,sec:ov-dynamics},
$\pi=\pi^J$ is the Gibbs measure at a fixed disorder. We write $\lambda$ for
the exponent of the barrier, and we use an upper bound
$\Delta_{n,\beta}\le C_bn^{-3\lambda}$ on the mean bias
$\Delta_{n,\beta}=P_\beta-n^{-1}\E F_n$. For \cref{thm:barrier-730}, we take
$\lambda=7/30$ and the bound $\Delta_{n,\beta}\le C_bn^{-7/10}$ of
\cref{inp:hb-rate}. The exponent of the upper bound is $1-\omega$, for a fixed
$\omega\in(0,1/4)$, and \cref{sec:ov-upper} uses the notation of
\cref{part:upper}.

\begin{table}[!t]
\centering
\scriptsize
\caption{The main results of
\cref{sec:model,sec:inputs,sec:scalar,sec:flow,sec:cubic,sec:fractional,sec:static,sec:references,sec:induction,sec:barrier-proof,u:sec:setup,u:sec:proposal,u:sec:fine,u:sec:concentration,u:sec:coarse,u:sec:connections,u:sec:proof,sec:fluct,sec:bias},
the sections whose results they use, and the sections in which they are used.
}\label{tab:ov-dependencies}
\begin{tabular}{@{}>{\raggedright\arraybackslash}p{0.17\textwidth}>{\raggedright\arraybackslash}p{0.5\textwidth}>{\raggedright\arraybackslash}p{0.1\textwidth}>{\raggedright\arraybackslash}p{0.13\textwidth}@{}}
\toprule
Section & Main results & Uses & Used in\\
\midrule
\ref{sec:model} Model & realized-start probability, \cref{lem:model-joint};
 heat-bath dynamics, \cref{lem:model-dynamics}; completed Hamiltonian,
 \cref{lem:model-completion} & --- & \ref{sec:inputs}, \ref{sec:fractional},
 \ref{sec:references}, \ref{sec:barrier-proof}, \ref{u:sec:setup},
 \ref{u:sec:connections}, \ref{u:sec:proof}, \ref{sec:fluct}, \ref{sec:bias}\\
\ref{sec:inputs} Inputs & matrix comparison, \cref{inp:ob-matrix}; upper
 bound on the mean bias, \cref{inp:hb-rate}; variance and profile bounds,
 \cref{inp:qp-rate,inp:qp-floor}; pair, fixed-scale, and separated-overlap bounds,
 \cref{inp:ob-pair-gap,inp:ob-fixed-gap,inp:ob-remote}; overlap-neighborhood
 masses, \cref{inp:ob-coverage} & \ref{sec:model}, literature &
 \ref{sec:scalar}--\ref{sec:references}, \ref{sec:barrier-proof},
 \ref{u:sec:setup}, \ref{u:sec:coarse}, \ref{u:sec:connections}\\
\ref{sec:scalar} Scalar problem & stationarity identities, \cref{b:lem:os-scalar-jets};
 the constant $m_\beta$, \cref{cor:sc-mbeta} & \ref{sec:inputs} & \ref{sec:cubic},
 \ref{sec:fractional}, \ref{sec:references}, \ref{u:sec:setup},
 \ref{u:sec:coarse}\\
\ref{sec:flow} Flows & uniform short-time bounds, \cref{prop:flow-uniform};
 response to the terminal function, \cref{b:lem:os-prefix,cor:flow-prefix} &
 \ref{sec:inputs} & \ref{sec:cubic}\\
\ref{sec:cubic} Cubic deficit & local expansion, \cref{b:prop:os-expansion};
 deficit on a band, \cref{prop:cub-band} & \ref{sec:inputs}--\ref{sec:flow} & \ref{sec:fractional}, \ref{sec:references}\\
\ref{sec:fractional} Fractional moments & fractional comparison,
 \cref{b:prop:fractional-matrix}; local fractional bound, \cref{lem:fr-local-gap}
 & \ref{sec:model}--\ref{sec:scalar}, \ref{sec:cubic} & \ref{sec:static}\\
\ref{sec:static} Static event & simultaneous disorder event, \cref{b:lem:static-event}
 & \ref{sec:inputs}, \ref{sec:fractional} & \ref{sec:references}, \ref{sec:barrier-proof}\\
\ref{sec:references} References & overlap-neighborhood masses,
 \cref{b:lem:coverage-extension}; fixed families of references,
 \cref{lem:ref-families} & \ref{sec:model}--\ref{sec:scalar}, \ref{sec:cubic}--\ref{sec:static} & \ref{sec:induction}, \ref{sec:barrier-proof}\\
\ref{sec:induction} Induction & constraints along a path, \cref{prop:ind-induction};
 a common time, \cref{cor:ind-witness} & \ref{sec:references} & \ref{sec:barrier-proof}\\
\ref{sec:barrier-proof} Barrier proof & sign obstruction, \cref{prop:bp-obstruction};
 \cref{thm:barrier-730,b:cor:subcritical,thm:barrier-strong} &
 \ref{sec:model}, \ref{sec:inputs}, \ref{sec:static}--\ref{sec:induction} &
 \ref{u:sec:proof}\\
\ref{u:sec:setup} Upper-bound setup & primitive disorder and Gaussian bounds,
 \cref{u:lem:goe-norm,u:lem:pressure-tail}; the scalar process beyond the
 support, \cref{u:prop:scalar,u:lem:cross,u:lem:tails} & \ref{sec:model},
 \ref{sec:inputs}, \ref{sec:scalar} & \ref{u:sec:proposal}--\ref{u:sec:proof}\\
\ref{u:sec:proposal} Proposal & Gaussian assembly, \cref{u:lem:assembly};
 spin laws, \cref{u:def:laws}; growing mesh, \cref{u:lem:quad} &
 \ref{u:sec:setup} & \ref{u:sec:fine}--\ref{u:sec:connections}\\
\ref{u:sec:fine} Fine entropy & reference law, \cref{u:lem:reference};
 forward entropy, \cref{u:lem:forward-kl,u:prop:fine-kl} &
 \ref{u:sec:setup}, \ref{u:sec:proposal} & \ref{u:sec:coarse},
 \ref{u:sec:connections}\\
\ref{u:sec:concentration} Concentration & fixed Gram stability,
 \cref{u:lem:gram}; weighted concentration, \cref{u:lem:weighted} &
 \ref{u:sec:setup}, \ref{u:sec:proposal} & \ref{u:sec:coarse}\\
\ref{u:sec:coarse} Coarse domination & weighted energy,
 \cref{u:lem:weighted-energy}; joint cancellation,
 \cref{u:lem:joint,u:lem:lag}; likelihood bound, \cref{u:prop:coarse} &
 \ref{sec:inputs}, \ref{sec:scalar}, \ref{u:sec:setup}--\ref{u:sec:concentration} &
 \ref{u:sec:connections}\\
\ref{u:sec:connections} Connections & smoothing and coverage,
 \cref{u:lem:smoothing,u:lem:coverage}; boundary mass, \cref{u:lem:boundary};
 connecting sequence, \cref{u:prop:sequence}; witness decay,
 \cref{u:prop:witness} & \ref{sec:model}, \ref{sec:inputs},
 \ref{u:sec:setup}--\ref{u:sec:fine}, \ref{u:sec:coarse}
 & \ref{u:sec:proof}\\
\ref{u:sec:proof} Upper-bound proof & pruning, \cref{u:lem:pruning};
 \cref{thm:upper,cor:upper-discrete,cor:two-sided}; parameters,
 \cref{u:prop:range} & \ref{sec:model}, \ref{sec:barrier-proof},
 \ref{u:sec:setup}, \ref{u:sec:connections} & ---\\
\ref{sec:fluct} Fluctuations & the bound of \cite{oai-fluctuations},
 \cref{f:window}; log-concave augmentation, \cref{f:logconcave}; \cref{f:main}
 & \ref{sec:model} & \ref{sec:bias}\\
\ref{sec:bias} Mean bias & variance--bias inequality,
 \cref{prop:bias-variance-deficit}; \cref{cor:mean-bias} & \ref{sec:model}, \ref{sec:fluct} & ---\\
\bottomrule
\end{tabular}
\end{table}

\section{The sign test}\label{sec:ov-sign}

Recall from \cref{lem:model-joint} that \cref{thm:barrier-730} states that
the Gibbs mass of the initial configurations from which the chain is at total
variation distance greater than $\tau$ from $\pi$ at time
$e^{A_\beta n^{7/30}}$ tends to one in probability over the disorder. Since $\pi$ is stationary, the chain started from a Gibbs sample has
law $\pi$ at every deterministic time (\cref{lem:model-dynamics}), and the law
averaged over the initial configuration carries no information. A lower bound
on the distance must come from a property of almost every individual initial
configuration that the dynamics forgets slowly.

The property we use is the sign of an overlap. As in
\cite[Section~8]{oai-barriers}, we choose for each $x\in\Sigma_n$ a reference
configuration $v(x)$, with $R(v(x),x)\ge q_*$ for all $x$ outside a set of small
Gibbs mass, where $q_*\in(0,1)$ is a fixed overlap (\cref{def:ind-algorithm}). Since $H_n(-y)=H_n(y)$, the half-space
$\{y:R(v(x),y)\le0\}$ has Gibbs mass at least $1/2$. If the law $\nu_x$ of the
chain started at $x$ is within $\tau$ of $\pi$ in total variation, then the
half-space has mass at least $1/2-\tau$ under $\nu_x$. Multiplying by
$\pi(x)$ and summing over $x$, we find that
\begin{equation}\label{eq:ov-half-space}
 \Bigl(\frac12-\tau\Bigr)\,\pi\bigl\{x:d^J_{\rm disc}(x,k)\le\tau\bigr\}
 \le\PP\bigl\{R(v(X_0),X_k)\le0\bigr\},
\end{equation}
where $(X_k)_{k\ge0}$ is the discrete chain with $X_0\sim\pi$
(see \eqref{eq:bp-half-space}). This is the reduction of
\cite[(8.1)--(8.3)]{oai-barriers}. The half-space depends on the
initial configuration, so that the distance on the left side is evaluated
before the average over $X_0$, while the right side is a probability for the
stationary chain. Since one update changes every overlap by at most $2/n$, a
chain that enters the half-space must cross every overlap level between $q_*$
and $0$.

We show that these crossings cannot occur along paths that avoid a fixed set
of small Gibbs mass. We call a sequence $x_0,\dots,x_m$ a path if all overlaps
change by at most $2/n$ from one element to the next \eqref{eq:ind-jump}, so
that every trajectory of the discrete chain is a path. Let $L=An^\lambda$ be
the logarithmic time scale, for a coefficient $A>0$ fixed below, so that $e^L$
is the time horizon, and let $T_n=\lceil n^2e^L\rceil$ be the budget of update
attempts. For every
disorder in an event $\mathcal G_n$ with $\PP(\mathcal G_n)\to1$, we construct
the map $x\mapsto v(x)$ and a set $V\subseteq\Sigma_n$ with
$\pi(V)\le3e^{-4L}$ such that
\begin{equation}\label{eq:ov-obstruction}
 R\bigl(v(x_0),x_m\bigr)>0\qquad\text{for every path }x_0,\dots,x_m
 \text{ with }m\le T_n\text{ that avoids }V
\end{equation}
(\cref{lem:ref-families,prop:bp-obstruction,lem:bp-coefficient}). Given
\eqref{eq:ov-obstruction}, a chain with $R(v(X_0),X_k)\le0$ for some
$k\le T_n$ visits $V$ at one of the times $0,\dots,k$. Since $X_a$ has law
$\pi$ for every deterministic $a$, a union bound implies that
\[
 \PP\bigl\{R(v(X_0),X_k)\le0\bigr\}\le(k+1)\pi(V)\le3(k+1)e^{-4L}
\]
(see \eqref{b:dyn:stationary-sign}). This is the only place where we use
stationarity, and we use it only at deterministic times. For this reason the
map $v$ and the set $V$ must be fixed before the chain is run, and they may not
depend on its trajectory. In continuous time, the number of updates up to time
$e^L$ is a Poisson variable with mean $ne^L$, and by Markov's inequality it
exceeds $T_n$ with probability at most $1/n$. We conclude that
\[
 \pi\bigl\{x:d^J(x,e^L)\le\tau\bigr\}\le\frac{3(T_n+1)e^{-4L}+1/n}{1/2-\tau}
\]
for every disorder in $\mathcal G_n$ (\cref{prop:bp-quenched}), and the right
side tends to zero. When we integrate over the disorder, $\PP(\mathcal G_n^c)$
is added once. The threshold $\tau$ enters only through the factor
$1/(1/2-\tau)$, which is why one coefficient serves every fixed $\tau<1/2$ and
both clocks.

\section{Overlap constraints and reference configurations}\label{sec:ov-architecture}

We follow the architecture of the proof of \cite[Theorem~1.1]{oai-barriers},
which proves a barrier at the time $\exp(n^{1/10000})$ and the threshold
$\tau=1/4$.

\paragraph{Overlap constraints.} For $(v,x,{\widetilde v})\in\Sigma_n^3$, we write
\[
 r=\frac{R(v,x)+R(v,{\widetilde v})}2,\qquad\delta=\frac{R(v,x)-R(v,{\widetilde v})}2,\qquad q=R(x,{\widetilde v}),
\]
so that the overlap matrix of $(v,x,{\widetilde v})$ is the matrix $D(r,q,\delta)$ of
\eqref{eq:inp-triple}. We use two kinds of constraints. The constraints of the
first kind act at fixed scales. For fixed cutoffs $0<b_2<b_1<1$
(\cref{tab:ref-parameters}), they state that if $R(x,{\widetilde v})\ge q_*$ and
$\abs{R(v,x)},\abs{R(v,{\widetilde v})}\le b_1$, then $\abs{R(v,x)-R(v,{\widetilde v})}\le2\eta_n$, and
similarly with the cutoffs $(b_1,b_2)$ in place of $(q_*,b_1)$. The triples
that violate them form the sets $\mathcal B_n(q_*,b_1)$ and
$\mathcal B_n(b_1,b_2)$ of \eqref{eq:ref-fixed-violation}. In words, a large
overlap between $x$ and ${\widetilde v}$ forces their small overlaps with
$v$ to agree. The
constraints of the second kind act at the scale $\rho$, near a level $r$ in a
compact interval $J_\beta\subset(0,q_\beta)$. With $\varepsilon=E\rho$ and
$E=1/100$, they state that no triple satisfies
\[
 r\in J_\beta,\qquad\varepsilon\le\abs\delta\le2\varepsilon,\qquad q\ge r+\rho.
\]
Such triples are called narrow, and they form the set $\mathcal N_n$
(\cref{sec:ref-parameters}). Equivalently, if $R(x,{\widetilde v})\ge r+\rho$, then
$\abs{R(v,x)-R(v,{\widetilde v})}$ does not lie in $[2\varepsilon,4\varepsilon]$. We
exclude a band, rather than a half-line, for the following reason. Along a
path, every overlap changes by at most $2/n<2\varepsilon$ in one step, and the
difference $\abs{R(v,x_k)-R(v,{\widetilde v})}$ cannot pass from below $2\varepsilon$ to
above $4\varepsilon$ without lying in the band at some time. As long as the path
keeps a large overlap with ${\widetilde v}$, a narrow constraint keeps its overlap with $v$
within $2\varepsilon$ of $R(v,{\widetilde v})$ (\cref{lem:ind-search}).

\paragraph{Constrained pressures.} In
\cref{sec:inputs,sec:scalar,sec:flow,sec:cubic,sec:fractional,sec:static}, we
write the three configurations as $(v,x,y)$. For a feasible overlap matrix
$D$, the probability that three independent samples from $\pi$ have overlap
matrix $D$ is $Z_{3,n}(D)/Z_n^3$, where $Z_{3,n}(D)$ is the partition function
restricted to the triples with overlap matrix $D$ (\cref{def:inp-constrained}).
There are at most $(n+1)^3$ feasible matrices, and it suffices to bound each
numerator from above and the denominator from below. Upper bounds for the
numerators come from Guerra's hierarchical interpolation~\cite{guerra-2003}
and its extension by Talagrand to replicas with prescribed overlaps
\cite[Section~3]{parisi-formula}, \cite[Propositions~3.4 and~4.4]{talagrand-2007}.
We use it in the form of \cite[Proposition~2.3]{oai-barriers}, stated as
\cref{inp:ob-matrix}. For every admissible path $(Q,m)$ of $3\times3$
covariance matrices, with positive semidefinite increments and nondecreasing
masses, the comparison with zero terminal multiplier bounds the constrained
pressure $p_{3,n}(D)=n^{-1}\E\log Z_{3,n}(D)$ by the trial value
\[
 p_{3,n}(D)\le\mathcal T_D(Q,m,0)
 =\Psi(Q,m,0)-\frac{\beta^2}{4}\int m\dd\norm Q_F^2
   +\frac{\beta^2}{4}\norm{D-Q_N}_F^2
\]
(\cref{def:inp-trial,inp:ob-matrix}). The last term is the cost of an endpoint
$Q_N$ that differs from $D$.

The baseline is the block hierarchy of
\cite[Section~4.2 and Section~7, Step~1]{oai-barriers}, described in
\cref{lem:inp-blocks} and built from the Parisi
distribution function $\alpha=\alpha_\beta$. Below the level $r$, the three
fields share their increments, at mass $\alpha/3$. Between $r$ and $q$, the
fields of $x$ and $y$ share their increments, at mass $\alpha/2$, while the
field of $v$ evolves alone, and above $q$ the three fields evolve
independently. The endpoint of this path is $D(r,q,0)$, and its trial value is
exactly $3P_\beta$. A perturbation of the baseline with endpoint
$D(r,q,\delta)$ and a smaller trial value proves a pressure deficit. Whether
such a perturbation exists is decided by the scalar Parisi problem. Let $u$ be
the spatial derivative of the solution of the Parisi equation, let $Y$ be the
associated diffusion, and let $\Gamma(t)=\E u(t,Y_t)^2$. It was proved in
\cite[Proposition~3.2]{oai-barriers} that $\Gamma(q)=q$ at every point $q$ of
the support of the Parisi measure, and by
\cite[Theorem~1.1]{full-rsb}, the support is $[0,q_\beta]$. We conclude that
$\Gamma(t)=t$ on $[0,q_\beta]$ (\cref{cor:inp-identity}). For two
configurations, we have the pair bound
$p_{2,n}(D_2(q))\le2P_\beta-(\Gamma(q)-q)^2/2$ of
\cite[Lemma~4.3]{oai-barriers} (\cref{inp:ob-pair-gap}). This deficit
vanishes on the support, and
constraints at overlaps inside $(0,q_\beta)$ require the three-replica
perturbations of \cite[Sections~4 and~7]{oai-barriers} and of
\cref{sec:cubic}.

\paragraph{Reference configurations and crossings.} In the dynamical argument,
$v$ and ${\widetilde v}$ are reference configurations and $x$ is a configuration of the
path. We take the references from lists of independent Gibbs samples, which
are fixed before any path is considered (\cref{sec:ref-families}). Suppose that
a path $x_0,\dots,x_m$ that avoids $V$ satisfies $R(v(x_0),x_m)\le0$. We
proceed through levels $i=0,1,\dots,N_{\rm lev}-1$, with thresholds
$r_0<r_1<\cdots$ in an interval $[t_0,2t_0]\subset J_\beta$
(\cref{def:ind-algorithm}). At level $0$, the reference is $v_0=v(x_0)$, and
$k_0$ is the first time with $R(v_0,x_k)\le r_0$. At level $j\ge1$, a bound
on the Gibbs mass of overlap neighborhoods supplies a new reference $v_j$ with
$R(v_j,x_{k_{j-1}})\ge q_*$ (\cref{b:lem:coverage-extension}, which extends
\cite[Proposition~6.1]{oai-barriers}). We then search backward in time from
$k_{j-1}$ for the last time $k_j$ with $R(v_j,x_k)\le r_j$. During the
search, the constraints keep the overlaps of the path with the earlier
references close to their values at $k_{j-1}$. The overlap
with $v_0$ stays close to $r_0<q_*$, while $R(v_0,x_0)\ge q_*$, so that the
search finds the crossing before it reaches time $0$ (\cref{lem:ind-crossing}).
After the last level, at the single time $k_{N_{\rm lev}-1}$, the overlaps
of the path with many references lie close to their thresholds. The
thresholds are drawn at random, and a count over the time index shows that this
is impossible when $m\le T_n$.

As in \cite[Sections~8.3 and~8.4]{oai-barriers}, the references are arranged
in blocks of $B$ levels. Within a block, a new reference has small overlaps with
the earlier references of the block, and it deletes every retained reference
with which its overlap exceeds $b_1$ in absolute value. Such conflicts occur
only at the first level of a block, and they never reach the first $B-H\ge2$
levels of the preceding block (\cref{lem:ind-deletion}).

\section{The exponent}\label{sec:ov-scales}

Three requirements fix the scales, which are listed in \cref{tab:ov-scales}.
The order in which the parameters are chosen is recorded in
\cref{tab:ref-parameters,tab:bp-order}.

\begin{table}[!t]
\centering
\small
\caption{The scales of the proof of \cref{thm:barrier-730} and their roles.
Here $R$, $E=1/100$, $S$, and $A_\beta$ are constants fixed before $n\to\infty$. The
last five rows are the rates of decay of the error terms of
\cref{b:lem:static-event}, without constants and polynomial factors; see
\cref{tab:st-errors,tab:ref-exponents}.}\label{tab:ov-scales}
\begin{tabular}{@{}ll>{\raggedright\arraybackslash}p{0.5\textwidth}@{}}
\toprule
Quantity & $\lambda=7/30$ & Role\\
\midrule
$\rho_n$ & $Rn^{-7/30}$ & gap $q-r$ of the narrow constraints;
 length of a threshold interval\\
$\varepsilon_n=E\rho_n$ & $ERn^{-7/30}$ & overlap imbalance
 $\varepsilon_n\le\abs\delta\le2\varepsilon_n$; radius of the windows\\
$\eta_n$ & $n^{-5/21}$ & tolerance of the fixed-scale constraints\\
$s_n=S\rho_n$ & $SRn^{-7/30}$ & order of the fractional moment\\
$L_n$ & $A_\beta n^{7/30}$ & logarithmic time scale; logarithm of
 the list sizes\\
\midrule
$n\rho_n^3$ & $R^3n^{3/10}$ & Gibbs probability of narrow triples\\
$n\eta_n^2$ & $n^{11/21}$ & Gibbs probability of fixed-scale triples\\
$n\rho_n^4$ & $R^4n^{1/15}$ & disorder failure of a local or
 separated-overlap bound\\
$n\eta_n^4$ & $n^{1/21}$ & disorder failure of a fixed-scale bound\\
$n^2\rho_n^6/n^{7/15}$ & $R^6n^{2/15}$ & disorder failure of the
 denominator (polynomial)\\
\bottomrule
\end{tabular}
\end{table}

\paragraph{The denominator.} On the disorder event of \cref{sec:static}, we
have $F_n\ge n(P_\beta-c_0\rho^3)$, so that
$Z_n^3\ge e^{3n(P_\beta-c_0\rho^3)}$. To obtain this bound from
Chebyshev's inequality with a fixed margin, we choose constants
$a_0<c_0$ such that $\Delta_{n,\beta}\le a_0\rho^3$
for all large $n$. Since $\E F_n=n(P_\beta-\Delta_{n,\beta})$, the distance from
the mean of $F_n$ to the threshold is then at least
$n(c_0-a_0)\rho^3$. The bounds on the numerators keep a deficit
$\varkappa g\rho^3$, where $g\rho^3$ is a lower bound, uniform over the band, for the
deficit per spin of the three-replica pressure near a level $r$
(\cref{sec:ov-cubic,sec:ov-static}) and $\varkappa\in(0,1)$ is fixed. The ratio
$Z_{3,n}(D)/Z_n^3$ is then at most $e^{-n(\varkappa g-3c_0)\rho^3}$, which is
small when $3c_0<\varkappa g$. Constants $a_0<c_0<\varkappa g/3$
with these properties exist if and only if
\[
 \limsup_{n\to\infty}\frac{\Delta_{n,\beta}}{\rho^3}<\frac{\varkappa g}3
\]
(see \eqref{eq:st-S1}). With $\Delta_{n,\beta}\le C_bn^{-7/10}$, this holds
for $\rho=Rn^{-7/30}$ if $R^3>3C_b/(\varkappa g)$ (see
\eqref{b:dyn:R-condition}). By Chebyshev's inequality and the bound
$\Var F_n\le C_vn^{7/15}$ of \cref{inp:qp-rate}, the lower bound on $F_n$ fails
with probability at most a constant times $n^{7/15}/(n^2\rho^6)$, which is
$O(n^{-2/15})$. We note that \cref{f:main} is not used in the proofs of
\cref{thm:barrier-730,thm:barrier-strong}.

\paragraph{The count.} The thresholds lie in the intervals
$I_i=[t_0+3i\rho,t_0+(3i+1)\rho]$, of length $\rho$ and spacing $3\rho$. There
are $N_{\rm lev}=MB$ levels, arranged in $M=\lfloor t_0/(3B\rho)\rfloor$
blocks (see \eqref{b:dyn:threshold-intervals}). In \cref{sec:ov-dynamics}, we
explain why a path of length $m$ that avoids $V$ and violates
\eqref{eq:ov-obstruction} produces, for every array of thresholds, a time
$k\le m$ at which $M$ overlaps lie in windows of radius
$\Psi_0=2\varepsilon+2\eta_n+2/n$ around their thresholds
(\cref{cor:ind-witness}). For independent thresholds, uniform on their
intervals, this event has probability at most
\begin{equation}\label{eq:ov-count}
 (m+1)\gamma_n^M,\qquad
 \gamma_n=\frac{2\Psi_0}\rho\longrightarrow4E=\frac1{25}
\end{equation}
(see \eqref{b:dyn:window-weights} and \eqref{b:dyn:prefix-count}). If this quantity
is less than one, such a path cannot exist. Note that $M$ is of order
$1/\rho$, so that we may take $\log m$ of order $n^\lambda$. This leads to the
time horizon $e^L$ with the logarithmic time scale $L=An^\lambda$, for every
$A<t_0\log25/(3BR)$ (see \eqref{b:dyn:coefficient}). The count requires the
ratio $\gamma_n$ of the length of a window to the length $\rho$ of a threshold
interval to stay below one. With $\varepsilon=E\rho$, the overlap
imbalance is of the largest permitted order, and this ratio stays below a
fixed bound. The gap satisfies $q-r\ge\rho$, and the weakest local deficit
occurs at the corner $q-r=\rho$, $\abs\delta=E\rho$ of the band of
\cref{sec:ov-cubic}.

In summary, the upper bound $n^{-3\lambda}$ on the mean bias allows
$\rho\asymp n^{-\lambda}$, there are then of order $n^\lambda$ levels, and the
count allows $\log t$ of order $n^\lambda$: the exponent of the barrier is one
third of the exponent of the upper bound on the mean bias. For the upper bound
$n^{-7/10}$ of \cref{inp:hb-rate}, this gives the exponent $7/30$ of
\cref{thm:barrier-730}.

\paragraph{The union bounds and the exponent $1/4$.} The proof uses two kinds
of union bound. In the disorder event, each of at most $(n+1)^3$ bounds on
constrained partition functions near a level $r$ fails with probability at
most $e^{-cn\rho^4}$ (\cref{sec:ov-static}). The sufficient condition
$n\rho^4/\log n\to\infty$ makes this polynomial union bound vanish (see
\eqref{eq:st-S3}). At a fixed disorder, we multiply a count $e^{CL}$ of reference
configurations and path positions by Gibbs probabilities bounded by
$C'(n+1)^3e^{-cn\rho^3}$, and it suffices that $cn\rho^3-CL-3\log(n+1)\to\infty$.
Under the stronger condition $n\rho^3/L\to\infty$, we have this for every fixed $C$.
For $\rho=Rn^{-\lambda}$ and $L=An^\lambda$, both sufficient conditions hold
when $\lambda<1/4$. At $\lambda=1/4$, the quantity $n\rho^4=R^4$ is bounded, and
the polynomial union bound in the disorder event is ineffective, whether or
not the conditional union bound can be absorbed by a choice of constants. The
argument gives no exponent of $1/4$ or more, whatever the upper bound on
the mean bias (\cref{tab:bp-exponents}). We note that a
mean bias of order $n^{-3/4}$, as predicted by Franz, Parisi, and
Ricci-Tersenghi~\cite{franz-parisi-ricci-2026}, would correspond exactly to this
threshold (\cref{sec:open}).

\section{The cubic deficit}\label{sec:ov-cubic}

We now describe the pressure deficit for the narrow constraints. We fix the
compact intervals $J_\beta\subset\operatorname{int}I_\beta$ with
$I_\beta\subset(0,q_\beta)$ of \cref{sec:ref-parameters}; the general
statements of \cref{sec:fractional,sec:static} write them as
$\mathcal J\subset\operatorname{int}I$. For a center $r\in J_\beta$, the gap
$d=q-r$, and the overlap imbalance $\delta$, we need a trial value for
$D(r,q,\delta)$ below $3P_\beta$ by a fixed multiple of $\rho^3$, uniformly
over $r\in J_\beta$ and over the band
$\{\rho\le d\le d_0,\ E\rho\le\abs\delta\le2E\rho\}$. The local cutoff $d_0$ is
chosen so small that every local window lies in $I_\beta$.

\paragraph{The obstacle.} In \cite[Section~7]{oai-barriers}, the field of $x$
receives $\delta$ times a Gaussian combination of the field increments in a
window $[r-h,r+h]$, and the field of $y$ receives $-\delta$ times the same
combination. The resulting excess covariance in the antisymmetric direction
$\mathbf e_x-\mathbf e_y$ contributes to the endpoint cost $(\beta^2/4)\norm{D-Q_N}_F^2$, and the expansion of
the trial value has the remainder $C(\abs\delta^3h^{-3/2}+\delta^4h^{-2})$
(see \eqref{eq:inp-remote}), which is small compared with the gain only if
$\abs\delta$ is much smaller than $h^{3/2}$, while $h$ is smaller than $q-r$.
In \cite{oai-barriers}, $h=\rho^8$ and $\abs\delta$ is of order $\rho^{20}$,
and the resulting deficit is of order $\rho^{43}$ \cite[(7.32)]{oai-barriers}.
When $\abs\delta$ and $q-r$ are both of order $\rho$, however, the mismatch and
the remainder are as large as the gain.

\paragraph{The observation.} For the perturbation constructed below, at
centers in the interior of the support, the terms of order at most two in the
width parameter $h$ of the window cancel exactly. The reason is the first stationarity identity
$\beta^2\E\chi(t,Y_t)^2=1$, where $\chi$ is the second spatial derivative of
the Parisi solution. This identity is the derivative of $\Gamma(t)=t$, and it
expresses the marginal stability of the Parisi solution
\cite[(22)--(23)]{crisanti-rizzo-2002}; its stochastic form is recorded in
\cite[Lemma~2.5]{full-rsb}. The deficit is then of
third order in $h$. To find it, we need a perturbed path whose endpoint
is exactly $D(r,q,\delta)$, so that we incur no mismatch cost, and an expansion of
its trial value to third order with a remainder of fourth order.

\paragraph{The perturbed path.} We fix a ratio $\vartheta\in(0,1/2]$; the
dynamical argument uses $\vartheta=1/2$. Given $(r,d,\delta)$, let $h=d/\vartheta$
be the width parameter of the window, $q_0=r-h$ its left endpoint, and $k=\delta/h$ the
relative imbalance (see \eqref{b:eq:os-parameters}). The perturbed path has the five
stages of \cref{tab:cub-stages}. In the baseline, the first four stages are
the following: the three fields share their increments on $[q_0,r]$, the
fields of $x$ and $y$ share theirs on $[r,q]$ while $v$ evolves alone, and all
three are independent on $[q,q_0+2h]$. To perturb them, we add $\delta\,\Xi$ to the
field of $x$ and subtract it from the field of $y$, where
\[
 \Xi=\frac1h\bigl\{\omega_0(g_1-g_2-g_x-g_y)+\omega_1g_v\bigr\},\qquad
 \omega_0=\frac1{2-\vartheta},\qquad\omega_1=\frac{1-\vartheta}{2-\vartheta},
\]
and $g_1,g_2,g_v,g_x,g_y$ are the field increments of the window
(see \eqref{eq:cub-contrast}). The weights are fixed, and the contrast extends past
the time $q$ at which $x$ and $y$ separate. By \eqref{eq:cub-contrast-moments},
the perturbation changes the covariance of the window in two ways: it shifts
the overlaps of $v$ with $x$ and with $y$ by $+\delta$ and $-\delta$, and it
adds an excess covariance of order $\delta^2/h$ in the antisymmetric
direction $\mathbf e_x-\mathbf e_y$. A fifth
stage, the terminal correction, removes this excess exactly. It lasts a time
proportional to $hk^2$, and it uses the positive semidefinite rate
$\mathrm{Id}-(\mathbf e_x-\mathbf e_y)(\mathbf e_x-\mathbf e_y)^{\mathsf T}/4$ in place of
$\mathrm{Id}$. Every increment of the path is positive semidefinite, the masses are
nondecreasing, and the endpoint is exactly $D(r,q,\delta)$
(\cref{lem:cub-identities,lem:cub-path}). From
\cite[Section~7]{oai-barriers} we keep the hierarchy in which the masses are
divided by the sizes of the common blocks, the antisymmetric Gaussian
contrast, and the source of the deficit in the identity
$\beta^2\E\chi(r,Y_r)^2=1$.

\paragraph{The reduction.} We split the trial value $\mathcal T(q_0,h,k)$
into the initial segment $[0,q_0]$, on which the three fields share every
increment, and the local window. Let $V(q_0,h,k,z)$ be the difference between
the value of the local window on the diagonal $\{z_v=z_x=z_y=z\}$ and its
unperturbed value $3\Phi(q_0,z)$ (see \eqref{eq:cub-V}), and let $\Delta\Pi$ be
the change of the penalty $\int m\dd\norm{Q}_F^2$. The initial segment responds
to the change $V$ of its terminal function linearly, up to an error
$\norm{V}_\infty^2/2$ that does not depend on the number of levels
(\cref{b:lem:os-prefix,cor:flow-prefix}). Since $\norm{V}_\infty\le Ch^2k^2$,
we have
\[
 \mathcal T(q_0,h,k)-3P_\beta=\E V(q_0,h,k,Y_{q_0})-\frac{\beta^2}4\Delta\Pi
 +O(h^4k^4)
\]
(\cref{cor:cub-prefix}). The expectation is taken under the law of the scalar
diffusion at the fixed time $q_0$, which does not depend on $(h,k)$; this is
why we expand at a fixed left endpoint. The local window is a composition of
five short semilinear parabolic flows on $\R^3$, some with singular diffusion
matrices. In \cref{sec:flow} we bound their mixed derivatives in $(h,k)$,
uniformly in the field and including $h=0$ (\cref{prop:flow-uniform}). Using
these bounds, we find that
\begin{equation}\label{eq:ov-pointwise}
 V(q_0,h,k,z)=\frac{\beta^4\chi_0^2}4\Delta\Pi
 +\beta^6h^3\bigl\{\alpha(q_0)^2\chi_0^3\,C_\chi(k)
 +\alpha(q_0)\,w_0^2\,C_w(k)\bigr\}+O(h^4k^2)
\end{equation}
(\cref{b:lem:os-flow,prop:cub-pointwise}). Here $\chi_0=\chi(q_0,z)$ and
$w_0=w(q_0,z)$ are the second and third spatial derivatives of the Parisi
solution, and $C_\chi$ and $C_w$ are the trace polynomial and the Wick
polynomial of \eqref{b:eq:os-R} and \eqref{b:eq:os-W}. When the mass is frozen
at $\alpha(q_0)$ on the window, the stages are Gaussian integration steps, and
three formal steps with explicit coefficient matrices give the coefficients of
order two and three (\cref{lem:cub-three-steps,lem:cub-centered,prop:cub-frozen}).
The variation of $\alpha$ on the window contributes, at order three, exactly
the corresponding part of $\beta^4\chi_0^2\Delta\Pi/4$ (\cref{lem:cub-mass}).

\paragraph{The cancellations.} The penalty enters the trial value twice:
directly, with the factor $-\beta^2/4$, and through $V$, with the factor
$\beta^4\chi(q_0,Y_{q_0})^2/4$. After the expectation, the sum is
$(\beta^2/4)(\beta^2\E\chi(q_0,Y_{q_0})^2-1)\Delta\Pi$, which vanishes by the
first stationarity identity. Since $\Delta\Pi$ is of order $h^2$,
this removes every term of order $h^2$, as well as the terms of order $h^3$ that involve
$\alpha'(q_0)$. The second stationarity identity
$\E w(t,Y_t)^2=2\alpha(t)\E\chi(t,Y_t)^3$ implies that the expectation of the
braces in \eqref{eq:ov-pointwise} equals
$\alpha(q_0)^2\E\chi(q_0,Y_{q_0})^3\{C_\chi(k)+2C_w(k)\}$. We conclude that
\begin{equation}\label{eq:ov-expansion}
 \Bigl|\mathcal T(q_0,h,k)-3P_\beta
 -\alpha(r)^2\beta^6\E\chi(r,Y_r)^3\,h^3P_\vartheta(k)\Bigr|\le Ch^4k^2,
 \qquad P_\vartheta=C_\chi+2C_w
\end{equation}
(\cref{b:prop:os-expansion}, with $q_0+h=r$). Both identities hold on
$(0,q_\beta)$ (\cref{b:lem:os-scalar-jets}), and both were known; see
\cite[proof of Lemma~7.2]{quantitative-parisi} and \cite[Lemma~2.5]{full-rsb}.
The use of the second identity in a covariance expansion is new. The
polynomial $P_\vartheta$ is even, with coefficients
$c_2(\vartheta)=2\vartheta(\vartheta-1)/(\vartheta-2)^2<0$,
$c_4(\vartheta)<0$, and $c_6(\vartheta)>0$ (see \eqref{eq:cub-c}). For
$\vartheta=1/2$,
\[
 P_{1/2}(k)=-\frac29k^2-\frac{62}{243}k^4+\frac{6952}{19683}k^6
\]
(see \eqref{eq:cub-P-half}), and \cref{sec:cub-coefficients} lists the traces from
which this polynomial can be recomputed by hand.

\paragraph{The deficit on the band.} In the variables $d$ and $\delta$, we
have $-h^3P_\vartheta(k)=\mathcal G_\vartheta(d,\abs\delta)$, where
\[
 \mathcal G_\vartheta(d,x)=-\frac{c_2}\vartheta\,x^2d-\vartheta c_4\frac{x^4}d
 -\vartheta^3c_6\frac{x^6}{d^3}
\]
(see \eqref{eq:cub-G-function}). The leading term is
$2(1-\vartheta)(2-\vartheta)^{-2}\delta^2d$, which is $4\delta^2d/9$ for
$\vartheta=1/2$. The deficit is quadratic in the overlap imbalance and
linear in the gap $q-r$; its minimum on the band is a positive multiple of
$\rho^3$, and it can be larger away from the corner. The scalar factor
$\alpha(r)^2\beta^6\E\chi(r,Y_r)^3$ is bounded below on $I_\beta$ by a constant
$m_\beta>0$, since $\alpha>0$ on $(0,q_\beta]$ and
$\E\chi(r,Y_r)^3\ge\beta^{-3}$ (\cref{cor:sc-mbeta}). For $E=1/100$, the
function $\mathcal G_\vartheta$ is increasing in $d$ and in $\abs\delta$ on the
band, and it attains its minimum at the corner $(\rho,E\rho)$
(\cref{lem:cub-band}). Choosing $d_0$ so small that the remainder in
\eqref{eq:ov-expansion} is absorbed, we obtain, for every fixed
$g'<m_\beta E^2G(\vartheta,E)$,
\begin{equation}\label{eq:ov-band}
 p_{3,n}\bigl(D(r,q,\delta)\bigr)\le\mathcal T(q_0,h,k)\le3P_\beta-g'\rho^3
\end{equation}
for every feasible matrix in the band, all $n$, and all $0<\rho\le d_0$,
uniformly over $r\in J_\beta$ (\cref{prop:cub-band}). Here $G(\vartheta,E)$ is the value of
$E^{-2}\rho^{-3}\mathcal G_\vartheta$ at the corner, and $G(1/2,1/100)>1/3$.
Since the second stationarity identity fails at $t=q_\beta$, where the Parisi
measure has an atom (\cref{rem:sc-upper-atom}), the centers lie in compact
subintervals of $(0,q_\beta)$; no other condition on the Parisi measure near
$r$ is needed, in contrast with \cite[Lemma~7.1]{oai-barriers}.

\section{From pressures to Gibbs probabilities}\label{sec:ov-static}

\paragraph{The obstacle.} The bound \eqref{eq:ov-band} concerns the quenched
pressure $n^{-1}\E\log Z_{3,n}(D)$. However, Gaussian concentration of this
Lipschitz function of the disorder (\cref{rem:inp-lipschitz}) turns a mean
deficit of order $\rho^3$ only into the probability bound $e^{-cn\rho^6}$, and
$n\rho^6$ tends to zero; for $\rho=Rn^{-7/30}$, it is $R^6n^{-2/5}$.

\paragraph{Fractional moments.} To overcome this obstacle, we bound positive
fractional moments of order $s=S\rho$, for a small constant $S>0$. Suppose that
\begin{equation}\label{eq:ov-fractional}
 \frac1{ns}\log\E Z_{3,n}(D)^s\le3P_\beta-g\rho^3 .
\end{equation}
Applying Markov's inequality to $Z_{3,n}(D)^s$, we find that
\[
 \PP\bigl\{n^{-1}\log Z_{3,n}(D)>3P_\beta-\varkappa g\rho^3\bigr\}
 \le e^{-(1-\varkappa)gS\,n\rho^4},
\]
and $n\rho^4$ grows like a power of $n$ (Step~2 of the proof of
\cref{b:lem:static-event}). To prove \eqref{eq:ov-fractional}, we adapt the
matrix comparison of \cite[Proposition~2.3]{oai-barriers} to moments of order
$s$, as in the positive-moment interpolation of Aronow and Lopatto
\cite[Section~2.3]{quantitative-parisi}: the disorder is integrated at the outer
mass $s$, instead of by an expectation. In this way we find that the trial
value of every admissible path whose masses are all at least $s$ bounds
$(ns)^{-1}\log\E Z_{\ell,n}(D)^s+\beta^2\ell^2s/(4n)$
(\cref{b:prop:fractional-matrix}). The interpolation is carried out for the
completed Hamiltonian $H_n+G$ of \cref{sec:model}, whose covariance is exactly
$nR(x,y)^2/2$. The correction $\beta^2\ell^2s/(4n)$ is quadratic in the number
$\ell$ of configurations, because the same variable $G$ enters every
configuration (\cref{lem:model-completion}). Since this correction is
nonnegative, it can be discarded.

\paragraph{The profile $\max\{\alpha,3s\}$.} The masses of the perturbed path
on its initial segment are $\alpha(t)/3$, and $\alpha(t)\to0$ as $t\to0$, so
that the requirement that all masses are at least $s$ fails near the origin.
To overcome this, we replace $\alpha$ by $\max\{\alpha,3s\}$ on the initial
segment. By the trial bound of Aronow and Lopatto
\cite[Lemma~7.2]{quantitative-parisi}, the Parisi functional of this profile
exceeds $P_\beta$ by at most $Cs^5$ (\cref{inp:qp-floor,lem:fr-floor}), which
is much smaller than $\rho^3$. The profile differs from $\alpha$ only on an
interval $[0,Cs]$, far from the window, which starts at $q_0\ge\min I_\beta>0$. We
show that its effect on the response of the initial segment to $V$ is at most
$C\norm{V'}_\infty\norm{\max\{\alpha,3s\}-\alpha}_{L^1}\le C\delta^2s^2/h$
(\cref{lem:fr-prefix-response}). On the band this is at most a constant times
$S^2\rho^3$, and we make it small by the choice of $S$
(see \eqref{b:eq:coarse-floor-error}). Together with \eqref{eq:ov-band}, this
proves \eqref{eq:ov-fractional} for every fixed $g<m_\beta E^2G(\vartheta,E)$,
for all $n$ and all small $\rho$, uniformly over the band
(\cref{lem:fr-local-gap}).

\paragraph{The other regimes.} Every narrow triple falls into one of three
regimes (\cref{lem:st-cover}). In the local regime $q-r\le d_0$, we use
\eqref{eq:ov-fractional}. In the regime $q-r>d_0$ of separated overlaps, with
$\abs{q-\Gamma(q)}$ small, we use the quadratic deficit $c_{\rm sep}\delta^2$ of
\cref{lem:st-remote}, which we derive from
\cite[Section~7, Steps~1--7]{oai-barriers} with fixed parameters; its source is
the decay $\beta^2\E[\chi(q,Y_q)\chi(r,Y_r)]\le1-c(q-r)$ of the mixed
susceptibility \cite[(7.31)]{oai-barriers}. Since $\delta^2$ is of order
$\rho^2$, Gaussian concentration suffices there. In the
remaining pair regime, the two configurations with the larger overlap $q$ have
$\abs{q-\Gamma(q)}>\nu$ for a fixed $\nu>0$, so that the pair deficit
$(\Gamma(q)-q)^2/2$ of \cite[Lemma~4.3]{oai-barriers} (\cref{lem:st-pair}) applies. For
the fixed-scale constraints, we use the three-replica deficit
$c_{\rm fix}\delta^2$ of \cite[(4.20)]{oai-barriers} (\cref{lem:st-fixed}) at
the tolerance $\abs\delta\ge\eta_n$. In the separated-overlap, pair, and
fixed-scale regimes, the mechanisms are those of \cite{oai-barriers}; what
changes are the scales at which the pressure deficits are converted into
probability bounds.

\paragraph{One disorder event.} The disorder event $\mathcal A_n$ of
\cref{b:lem:static-event} requires the lower bound on $F_n$ of
\cref{sec:ov-scales} and upper bounds on all constrained partition functions
of the local, separated-overlap, fixed-scale, and pair families; these are
deterministic families of at most $C(n+1)^3$ matrices. The event
$\mathcal A_n$ fails with probability at most
\[
 C\frac{n^{7/15}}{n^2\rho^6}+C(n+1)^3\bigl(e^{-cn\rho^4}+e^{-cn\eta_n^4}\bigr)
 +C(n+1)e^{-cn}
\]
(see \eqref{b:eq:static-event-failure}), which is $O(n^{-2/15})$ for
$\lambda=7/30$. For every disorder in $\mathcal A_n$,
dividing the numerators by the denominator, we find that
\[
 \pi^{\otimes3}(\mathcal N_n)\le C(n+1)^3e^{-cn\rho^3},\qquad
 \pi^{\otimes3}(\mathcal W_n)\le C(n+1)^3e^{-cn\eta_n^2},
\]
where $\mathcal W_n$ is the set of fixed-scale triples
(see \eqref{b:eq:conditional-static-errors}). We keep two kinds of error apart
(\cref{rem:st-two-levels}). Our estimate for the disorder failure probability
$\PP(\mathcal A_n^c)$ contains a polynomial term, and we use this estimate
once, additively, at the end of the proof. The Gibbs
probabilities at a fixed disorder in $\mathcal A_n$ decay at least as fast as
$e^{-cn^{1-3\lambda}}$, and only these are multiplied by the numbers $e^{O(L)}$
of reference configurations and path positions. This separation replaces the
conversion \cite[Lemma~4.2]{oai-barriers} of pressure bounds into probability
bounds, whose finite-size error $n^{-1/24}$ is far too large at the scale
$\rho^3$.

\section{The dynamical argument}\label{sec:ov-dynamics}

\paragraph{Fixed families of references.} The dynamical argument starts from
two inputs at a fixed disorder. The first is the bound on the Gibbs mass of
overlap neighborhoods: for every $S\subseteq\Sigma_n$, all $x\in S$ outside a
set of Gibbs mass less than $e^{-D_1L}$ have Gibbs mass at least $e^{-D_3L}$ of
configurations $z\in S$ with $\abs{R(x,z)}\ge q_*$. We prove it by the argument
of \cite[Proposition~6.1]{oai-barriers}, at every scale $L\to\infty$ with
$L=o(n)$ (\cref{b:lem:coverage-extension}). The second input is the bound
$\pi^{\otimes3}(\mathcal B_n)\le q_n$ for the exceptional triples
$\mathcal B_n=\mathcal N_n\cup\mathcal B_n(q_*,b_1)\cup\mathcal B_n(b_1,b_2)$ on
the event $\mathcal A_n$, where
$q_n=C(n+1)^3(e^{-cn\rho^3}+e^{-cn\eta_n^2}+e^{-cn})$ satisfies
$e^{CL}q_n\to0$ for every fixed $C$ (\cref{lem:ref-single-event}).

In \cite[Section~8.3]{oai-barriers}, the properties of the lists of references
are proved with probability tending to one jointly over the disorder, the
lists, and the path, with bounds averaged over the disorder. At our scales,
such an averaged bound contains the disorder failure probability, which does
not remain small after multiplication by $e^{O(L)}$. For this reason, we fix
the disorder first. At each
disorder in $\mathcal G_n$, we sample the lists, with
$\lceil e^{(c_s+1)L}\rceil$ independent Gibbs samples at the levels of
position $s$ in a block. The expected Gibbs mass of the configurations at which
some cell contains no suitable entry is at most $2e^{-4L}$. The expected number
of exceptional triples of list entries, and the expected Gibbs mass of the
configurations that form one with two entries, are at most $e^{O(L)}q_n$. We fix
a realization at which the sum of such quantities is at most its mean
(\cref{lem:ref-families}). Outside a set $V$ with $\pi(V)\le3e^{-4L}$, it has
properties (P1)--(P4) of \cref{sec:ind-setting}: every cell supplies a new
reference for each configuration in it, the fixed-scale constraints hold across $q_*$ and across $b_1$, and
no configuration outside $V$ forms a narrow triple with two signed entries.
The family and $V$ depend only on $J$ (\cref{rem:ref-every-path}).

\paragraph{The accumulation of errors.} In
\cite[(8.19)--(8.20)]{oai-barriers}, the overlaps of every retained reference
with the current configuration and with the later references are controlled
within $6j\varepsilon$ after $j$ levels. These errors accumulate, and the final
count of \cite[(8.29)]{oai-barriers} requires $N_{\rm lev}\varepsilon$ to be
much smaller than $\rho$. Since $N_{\rm lev}$ is of order $1/\rho$, this forces
$\varepsilon$ to be much smaller than $\rho^2$. Our choice $\varepsilon=E\rho$,
at which the narrow constraints have a deficit of order $\rho^3$, is far outside
this range.

\paragraph{Next retained references.} To overcome this, we compare each
retained reference $v_i$ with the next retained reference $v_{a_j(i)}$, where
$a_j(i)$ is the smallest retained level larger than $i$ after level $j$
(\cref{sec:ind-algorithm}). We call the overlap $c_j(i)=R(v_i,v_{a_j(i)})$ the
center of $v_i$. When $v_{i+1}$ is selected at the time $k_i$, we have
$R(v_{i+1},x_{k_i})\ge q_*$ and $R(v_i,x_{k_i})\in[r_i-2/n,r_i]$, and the
fixed-scale constraint across $q_*$ implies that
$\abs{c_j(i)-r_i}\le2\eta_n+2/n$. The center changes only when a deletion
removes the next retained reference of $i$. The next retained reference is then
the first reference of the following block, which is never deleted, and the
center moves by at most $2\varepsilon+2\eta_n$. We show by induction that, for
every retained $i<j$,
\[
 \abs{c_j(i)-r_i}\le2\varepsilon+4\eta_n+\frac2n,\qquad
 \abs{R(v_i,x_{k_j})-c_j(i)}<2\varepsilon
\]
(\cref{def:ind-invariant,prop:ind-induction}). The errors are of order
$\varepsilon$ at every level, and they do not grow with the number of levels.

The narrow constraints preserve these windows during the backward search
(\cref{lem:ind-search}). Suppose that some window fails at a time between
$k_j$ and $k_{j-1}$, and let $k$ be the largest such time and $v_i$ a
reference whose window fails at $k$. All windows hold at time $k+1$, and since
one step changes overlaps by at most $2/n$, the overlap $R(v_i,x_k)$ differs
from $c_j(i)$ by an amount in $[2\varepsilon,2\varepsilon+2/n)$. Let
$a=a_j(i)$. Then $R(v_a,x_k)$ is close to $r_a$ or larger: if $a=j$, because the
backward search has not yet reached the crossing of $r_j$, and otherwise
because the window of $v_a$ holds at time $k+1$. Since $r_a-r_i\ge2\rho$ and
all errors are of order $\varepsilon$, we find that $R(v_a,x_k)\ge r+\rho$, where $r$ is the mean of $R(v_i,x_k)$ and
$R(v_i,v_a)=c_j(i)$. The triple $(v_i,x_k,v_a)$ is then narrow, which (P4)
excludes because $x_k\notin V$. In other words, the next retained reference
supplies the large overlap that the narrow constraint needs, and the spacing
$3\rho$ of the threshold intervals makes this overlap large enough.

\paragraph{The count.} At the time $k_*=k_{N_{\rm lev}-1}$, the first
reference $v_{\ell B}$ of each block $\ell$ still has its original next retained
reference $v_{\ell B+1}$, because the first two levels of a block are never
deleted. By its window, we have $\abs{R(v_{\ell B},x_{k_*})-r_{\ell B}}<\Psi_0$ for
$\ell=0,\dots,M-1$ (\cref{cor:ind-witness}). The reference of each level depends
only on the thresholds of the earlier levels (\cref{lem:ind-predictable}); this
is why a selected reference whose crossing fails is kept. Removing the windows
one at a time, from the largest level down, we find that for each fixed time
$k$ the $M$ windows hold with probability at most $\gamma_n^M$ over independent
uniform thresholds (\cref{lem:bp-product}). Summing over $k\le m$, we have the
bound \eqref{eq:ov-count}. If $(m+1)\gamma_n^M<1$, some threshold array admits no such
time, which contradicts the existence of a path that violates
\eqref{eq:ov-obstruction} (\cref{prop:bp-obstruction}). In \cite[(8.29)]{oai-barriers}, the count is a union over
subsets of retained levels; here the windows sit at the deterministic first
levels of the blocks, and we count only $M$ of them. Further, since the
static estimates are uniform over the interval $J_\beta$, we need no random
event that restricts the thresholds, such as \cite[(8.12)]{oai-barriers}.

\paragraph{Quantitative and almost-sure forms.} With the lower bound on $F_n$
from the small-deviation inequality of Paouris and Valettas
\cite[Theorem~2.1 and Remark~2.4(3)]{paouris-valettas-2018} instead of
Chebyshev's inequality, with $\eta_n=\rho_n/1000$, and with
$A<t_0\log(250/11)/(3BR)$, the disorder event fails with probability at most
$Ce^{-cn^{1/15}}$, and on it the initial configurations from which the
distance stays at least $1/2-Ce^{-cn^{7/30}}$ up to the time horizon
$e^{An^{7/30}}$ have Gibbs mass at least $1-Ce^{-cn^{7/30}}$, for all
$n\ge n_1(C_{\rm rate})$ and every single-site generator with rates at most a
fixed $C_{\rm rate}$ that leaves $\pi$ stationary (\cref{thm:barrier-strong}).
By the same theorem, we have the relaxation time bound
$e^{An^{7/30}}/(Cn^{7/30})$ for the reversible ones and, by the
Borel--Cantelli lemma, the almost sure form for every coupling of the
disorders of different sizes.

\section{The upper bound}\label{sec:ov-upper}

We follow the architecture of the proof of
\cite[Theorem~1]{oai-typical-upper}, which proves \cref{thm:upper} for the
saving $\omega=1/40000000$. The proof has four steps. A finite-chain argument
reduces a bound for typical initial configurations to a lower bound on the
stationary flow across every cut whose two sides both have appreciable Gibbs
mass. A random disorder and a random configuration are then constructed
jointly, by Gaussian queries of the disorder matrix that follow the scalar
Parisi diffusion. The law of this pair is compared with the Gibbs joint law in
two ways: through its relative entropy, and through an upper bound on its
likelihood ratio outside a rare event. A Gaussian smoothing of the disorder
reverses the entropy comparison, so that the constructed configuration reaches
every set of appreciable Gibbs mass. Finally, a sequence of such
configurations, built from rotated Gaussian inputs, connects the two sides of
a cut of small flow while avoiding its boundary, which is impossible. In
\cite{oai-typical-upper}, the parameters are fixed numbers, and the exponent
is not optimized.

To reach every $\omega<1/4$, two estimates must be close to sharp: the
relative entropy of the proposal at the longer horizon, and the likelihood
bound at the shorter horizon, together with the probability of its
exceptional event. The first is described in the paragraph on the forward
entropy below, and the second in the paragraph on coarse domination. The other
steps follow \cite{oai-typical-upper}, with the changes of parameters that
these two estimates allow.

\paragraph{Notation.} Let $g\in\R^{n\times n}$ have independent standard
Gaussian entries, with law $\gamma_n$, and let
$W(g)=(g+g^{\mathsf T})/\sqrt{2n}$. With $J_{ij}=(g_{ij}+g_{ji})/\sqrt2$ and
$J'_{ii}=g_{ii}$, the matrix $W(g)$ is the matrix $W$ of
\cref{sec:model-completion}, and $H^c_g(x)=x^{\mathsf T}W(g)x/2$ is the
completed Hamiltonian of the disorder $J(g)$. We write $F^c(g)$ for the
completed log partition function, $\pi_g=\pi^{J(g)}$, $P^g_t=P^{J(g)}_t$, and
\[
 \Pi_n(\dd g,x)=\gamma_n(\dd g)\,\pi_g(x)
\]
for the Gibbs joint law. Its $(J(g),x)$ marginal is the realized-start law
$\PP_{J,X_0\sim\pi^J}$ of \cref{def:model-joint}. For $S\subseteq\Sigma_n$,
the stationary flow across $S$ is
\[
 \mathcal Q_g(S,S^c)=\sum_{x\in S,\,y\notin S}\pi_g(x)\,\mathcal L^{J(g)}(x,y),
\]
where $\mathcal L^{J(g)}(x,x^{i,-x_i})=p_i^{J(g)}(-x_i\mid x)$ is the rate of
a flip at site $i$ (\cref{sec:model-dynamics}). The exponents of the proof
are fixed numbers, written in sans-serif ($\ex a,\ex q,\ex k_{\rm f},\dots$);
they are chosen after $\beta$ and $\omega$, and they are listed in
\cref{tab:ov-upper-exponents}. Constants $C_\beta,c_\beta>0$ depend only on
$\beta$ and on the exponents.

\begin{table}[!t]
\centering
\small
\caption{The exponents of the proof of \cref{thm:upper}, as functions of
$\varepsilon=\varepsilon_\omega=(1/4-\omega)/20$, their values for
$\omega=1/5$, and their roles. They are fixed after $\beta$ and $\omega$, and
before $n\to\infty$; see \cref{u:sec:parameters,u:sec:budget}.}
\label{tab:ov-upper-exponents}
\begin{tabular}{@{}lll>{\raggedright\arraybackslash}p{0.5\textwidth}@{}}
\toprule
Exponent & Value & $\omega=1/5$ & Role\\
\midrule
$\ex a$ & $1/2$ & $1/2$ & mesh $h=n^{-\ex a}$\\
$\ex f$ & $1/2$ & $1/2$ & fine horizon $V_{\ex f}\approx\ex f\log n$\\
$\ex b$ & $1/2-14\varepsilon$ & $93/200$ & coarse horizon $V_{\ex b}\approx\ex b\log n$\\
$\ex k_{\rm f}$ & $1/2-2\varepsilon$ & $99/200$ & forward entropy $n^{1-\ex k_{\rm f}}$\\
$\ex D$ & $1/4-\varepsilon$ & $99/400$ & smoothing $\delta=n^{-\ex D}$\\
$\ex u_{\rm E}$ & $1/4-2\varepsilon$ & $49/200$ & reverse entropy $n^{1-\ex u_{\rm E}}$\\
$\ex d$ & $\varepsilon/2$ & $1/800$ & masses $n^{-\ex d}$ of the cuts\\
$\ex e$ & $1/4-3\varepsilon$ & $97/400$ & crossing probability $e^{-2n^{1-\ex e}}$\\
$\ex w$ & $1/8-3\varepsilon$ & $47/400$ & weighted errors $n^{-\ex w}$\\
$\ex z$ & $1/4-8\varepsilon$ & $23/100$ & accuracy $n^{-\ex z}$ of the row averages\\
$\ex k_{\rm b}$ & $1/4-9\varepsilon$ & $91/400$ & likelihood bound $e^{n^{1-\ex k_{\rm b}}}$\\
$\ex s_{\rm b}$ & $1/4-5\varepsilon$ & $19/80$ & failure $e^{-n^{1-\ex s_{\rm b}}}$ of the coarse proposal\\
$\ex s_{\rm seq}$ & $1/4-4\varepsilon$ & $6/25$ & failure $e^{-n^{1-\ex s_{\rm seq}}}$ of the sequence\\
$\ex r$ & $1/4-15\varepsilon/2$ & $37/160$ & Hamming radius $R_n=\lceil n^{1-\ex r}\rceil$\\
$\ex u_{\rm F}$ & $1/4-39\varepsilon/4$ & $361/1600$ & threshold $n^{1-\ex u_{\rm F}}$ for sums of local fields\\
$\ex q$ & $1/4-10\varepsilon$ & $9/40$ & flow threshold $\phi_n=e^{-n^{1-\ex q}}$\\
$\ex z_{\rm P}$ & $3/10$ & $3/10$ & pressure event $F^c\le nP_\beta+n^{1-\ex z_{\rm P}}$\\
$\ex g_{\rm rot}$ & $3/5$ & $3/5$ & grid of $\ell_n+1\approx e^{n^{\ex g_{\rm rot}}}$ angles\\
\bottomrule
\end{tabular}
\end{table}

\paragraph{From balanced-cut flow to mixing.} By \cref{lem:model-joint},
\cref{thm:upper} states that the Gibbs mass of the configurations $x$ with
$d^{J(g)}(x,e^{n^{1-\omega}})>1/4$ tends to zero in probability over $g$. The
reduction to flows is the finite-chain lemma
\cite[Lemma~8]{oai-typical-upper}, which we prove as \cref{u:lem:pruning}.
Let $0<\epsilon<1/3$ and $\phi>0$, and suppose that
$\mathcal Q_g(S,S^c)>\phi$ for every set $S$ with
$\epsilon\le\pi_g(S)\le1-\epsilon$. Then, for every $t>0$,
\begin{equation}\label{eq:ov-pruning}
 \sum_{x\in\Sigma_n}\pi_g(x)\,d^{J(g)}(x,t)
 \le\frac{n\sqrt{2\log2}}{\phi\sqrt t}+2\epsilon .
\end{equation}
Here $n$ bounds the total jump rate out of a configuration, and $n\log2$
bounds the entropy of $\pi_g$. The proof removes, one at a time, sets of at
most half of the remaining mass whose flow to the rest is less than $\phi$
times their mass. The hypothesis prevents the removed mass from reaching $\epsilon$.
On the remaining set, the flow bound gives a Poincar\'e inequality, by the
argument of Lawler and Sokal \cite[Theorem~3.5]{lawler-sokal-1988}, and the
entropy dissipation of the full chain
\cite[Lemmas~2.5 and~2.7]{diaconis-saloff-coste-1996} bounds the average of
the distances from fixed initial configurations. The trajectory need not
stay in the remaining set. Initial configurations from which the chain is
slow are allowed, and a set of them costs at most its Gibbs mass. For this
reason only cuts of mass at least $\epsilon$ enter, and the conclusion
concerns typical initial configurations; at large $\beta$, the worst-case
mixing time is exponential \cite[Theorem~1.1]{sellke-2025}.

We apply \eqref{eq:ov-pruning} with $\epsilon=n^{-\ex d}$,
$\phi=\phi_n=e^{-n^{1-\ex q}}$, and $t=e^{n^{1-\omega}}$. By Markov's inequality
in the Gibbs draw, we find that the joint probability that
$d^{J(g)}(x,t)>1/4$ is at most
\[
 \gamma_n(\mathcal W_n)+2e^{-n}
 +4n\sqrt{2\log2}\,e^{n^{1-\ex q}-n^{1-\omega}/2}+8n^{-\ex d}
\]
(\cref{u:sec:proof}). Since $\omega<\ex q$, the third term tends to zero.
The term $2e^{-n}$ bounds the probability that
$\norm{W(g)}_{\rm op}>12$ (\cref{u:lem:goe-norm}), and
\[
 \begin{split}
 \mathcal W_n=\bigl\{g:{}&\norm{W(g)}_{\rm op}\le12,\text{ and there exists }S
 \subseteq\Sigma_n\\
 &\text{with }n^{-\ex d}\le\pi_g(S)\le1-n^{-\ex d}\text{ and }
 \mathcal Q_g(S,S^c)\le\phi_n\bigr\}
 \end{split}
\]
is the event that a cut of small flow exists. The rest of the proof shows
that $\gamma_n(\mathcal W_n)\to0$ (\cref{u:prop:witness}).

\paragraph{The proposal.} The flow bound is proved with a random pair $(g,x)$
whose law is close to $\Pi_n$. Its construction is that of
\cite[Section~4]{oai-typical-upper}, which is related to the iterative
construction of solutions of the TAP equations by Bolthausen
\cite{bolthausen-2014} and to the incremental approximate message passing algorithm of
Montanari \cite{montanari-2021}. Let $\alpha_\beta$ be the Parisi distribution
function; the Parisi measure has support $[0,q_\beta]$ (\cref{inp:full-rsb}).
In the variance
time $v=\beta^2t$, the scalar diffusion of \cref{sec:inputs} solves
\[
 \dd Y_v=\dd M_v+\bar\alpha(v)\,m(v,Y_v)\dd v,\qquad Y_0=0,
\]
where $M$ is a standard Brownian motion, $\bar\alpha(v)=\alpha_\beta(v/\beta^2)$,
and $m(v,y)=u_{\alpha_\beta}(v/\beta^2,y)$ is the spatial derivative of the
Parisi solution. The magnetization $m_v=m(v,Y_v)$ is a martingale, with
$\dd m_v=\chi_v\dd M_v$, where $\chi_v=\chi(v,Y_v)$ is the second spatial
derivative. Beyond $v_*=\beta^2q_\beta$, we continue the process with
$\bar\alpha=1$, $m(v,y)=\tanh y$, and $\chi(v,y)={\rm sech}^2y$; the continued
potential solves the same equation (\cref{u:sec:scalar}). We define
\[
 p(v)=\E\chi_v,\qquad r(v)=\E\chi_v^2,\qquad
 \rho(v)=\beta\sqrt{r(v)},\qquad \sigma(v)=\sqrt{1-\rho(v)^2}.
\]
By \cref{cor:inp-identity}, we have $r=\beta^{-2}$ on $[0,v_*]$, so that
$\rho=1$ and $\sigma=0$ there. After $v_*$, the functions $p$ and $r$ decay at
the rate $e^{-v/2}$ up to a factor $\sqrt{1+v}$
(\cref{u:prop:scalar,u:lem:tails}), and at a horizon $V$ of order $\log n$ the
magnetization is close to $\pm1$.

In the proposal, we run one copy of the continued scalar process for each
coordinate, driven by noise extracted from the matrix $W=W(g)$. We cut the time
axis into cells $[t_j,t_{j+1}]$ with
\[
 t_0=0,\qquad h_j=he^{t_j/4},\qquad t_{j+1}=t_j+h_j,\qquad h=n^{-\ex a}.
\]
At the query $j\ge1$, we take the increment of the vector of magnetizations
over the previous cell, normalized as
$\mathcal T_j=(m(t_j)-m(t_{j-1}))/s_j$ with $s_j^2=\int_{t_{j-1}}^{t_j}r$, and
we orthogonalize it against the earlier query directions in the inner product
$n^{-1}\langle u,v\rangle$. This gives a direction $u_j$, and we read the
column $Wu_j$. Its component orthogonal to $u_1,\dots,u_j$, together with an
independent Gaussian completion inside their span, is a vector $e_j$ that is
standard Gaussian given the past. The vector $e_j$ becomes the standardized
increment of the query paths $B_i$ over the next cell $[t_j,t_{j+1}]$, which
is filled with independent Brownian bridges. Each coordinate is driven by
$M_i=\int\rho\dd B_i+\int\sigma\dd D_i$, where the $D_i$ are independent
Brownian motions. Up to $v_*$, the matrix drives the process alone; after
$v_*$, the query noise is mixed with independent noise. By Gaussian orthogonal
invariance, the unread part of $W$ is a compressed Gaussian orthogonal
ensemble, independent of the past. Then the matrix assembled from the answers
and the remaining Gaussian variables has the law of $W(g)$ under $\gamma_n$,
independently of the initial data, and the pairs $(B_i,D_i)$ are independent
pairs of Brownian motions (\cref{u:lem:assembly}). At the horizon $V$, the
spin is linked to the terminal magnetizations by the density
$\prod_i(1+x_im_i(V))$ relative to an independent uniform spin
(\cref{u:def:laws}). The output $Q_V$ is the law of $(g,x)$, where $g$ is
obtained from $W$ by adding an independent antisymmetric Gaussian matrix, so
that $W(g)=W$. Its marginals are
$\gamma_n$ and the uniform law, but $x$ is correlated with $g$.

The construction has two equivalent forms: in one, the matrix is assembled
from the answers to the queries, and in the other, the same recursion reads a
given matrix and given completions (\cref{u:def:laws,u:lem:assembly}). The main
structural fact is that, under the linked law, the rows $(x_i,Y_i,B_i,D_i)$,
$i=1,\dots,n$, are independent, with an explicit one-row law $\nu_V$.
Apart from the Gaussian variables that the recursion never reads, the
statistics of the proposal used below are averages of $n$ independent rows, or
explicit functions of such averages.

We use two horizons on the same mesh: the fine horizon $V_{\ex f}$ and the
coarse horizon $V_{\ex b}$, the first times $t_j$ beyond $\ex f\log n$ and
$\ex b\log n$, with $\ex b<\ex f$ (\cref{u:lem:quad}). The two runs read the
same inputs, so that the coarse run is an initial segment of the fine run. The
growing cells are new compared with \cite{oai-typical-upper}, where the cells
have equal length $n^{-1/16}$. Since $r(t_j)$ decays like $e^{-t_j/2}$, the
lengths $h_j$ grow as the profile decays, and $h_j^2r(t_j)$ stays of order
$h^2$ at most. The number of queries is at most $8n^{\ex a}$ for every horizon
(\cref{u:lem:quad}), and the errors of the mesh add up to $O(h)$, uniformly in
the horizon. Unequal cells require some care in the estimates: $\mathcal T_j$
is an increment over the previous cell, $e_j$ belongs to the next cell, and the
ratio $h_j/h_{j-1}=e^{h_{j-1}/4}$ enters the bounds.

\paragraph{The forward entropy.} The first comparison is
\begin{equation}\label{eq:ov-forward-kl}
 \KL(Q_V\Vert\Pi_n)\le nj(V)-n\log2+\E F^c(g)-\beta\,\E H^c_g(x),
 \qquad j(V)=\frac12\int_0^V\rho(v)^2\varkappa(v)\dd v,
\end{equation}
where $\varkappa(v)=\E[\chi_v^2/(1-m_v^2)]$ and the last expectation is under
the proposal (\cref{u:lem:forward-kl}). The quantity $j(V)$ is the
information per spin that the proposal acquires through the matrix. To see
this, we move the part of the linking density that comes from the independent
noise $D$ into a reference law, as in \cite[Section~4]{oai-typical-upper}, whose $(g,x)$ marginal is $\gamma_n$ times the
uniform law. The logarithm of the remaining density is a sum of $n$
independent rows, each with mean $j(V)$ (\cref{u:lem:reference}). The factor
$\rho^2$ is the share of the matrix in the driving noise. Changing the
reference to $\Pi_n$, whose density is $2^ne^{\beta H^c_g(x)-F^c(g)}$, gives
\eqref{eq:ov-forward-kl}.

The terms of \eqref{eq:ov-forward-kl} are related by the scalar identity
\cite[Lemma~4]{oai-typical-upper}: for every $V\ge v_*$,
\begin{equation}\label{eq:ov-scalar-identity}
 j(V)-\log2+P_\beta-\beta e(V)=\frac{\beta^2}4p(V)^2,
 \qquad e(V)=\beta\int_0^Vr(v)p(v)\dd v
\end{equation}
(\cref{u:prop:scalar}). Here $e(V)$ is the scalar reference for the energy per spin of the proposal,
and the identity balances the information paid through the matrix against the
energy gained. By Guerra's bound and \cref{lem:model-completion}, we have
$\E F^c(g)\le nP_\beta$ (\cref{inp:parisi}). The expected energy satisfies
\[
 \frac1n\E H^c_g(x)\ge e(V)-C_\beta\Bigl[h+\frac{(N+1)(1+V)+\sqrt N}n
 +\frac{(N+1)^2VR_V}{n^2}\Bigr],
\]
where $N$ is the number of queries and $R_V=(\min_{[0,V]}r)^{-1}$
(\cref{u:lem:fine-energy}). Inserting these bounds into
\eqref{eq:ov-forward-kl}, and using $p(V)^2\le C_\beta e^{-V}$, we find that
\[
 \KL(Q_V\Vert\Pi_n)\le C_\beta\Bigl[nh+ne^{-V}+(N+1)(1+V)+\sqrt N
 +\frac{(N+1)^2VR_V}n\Bigr].
\]
At the fine horizon, the five terms have the powers $1-\ex a$, $1-\ex f$,
$\ex a$, $\ex a/2$, and $2\ex a+\ex f/2-1$ of $n$, up to logarithms, and we
obtain $\KL(Q_F\Vert\Pi_n)\le C_\beta n^{1-\ex k_{\rm f}}$ for the fine
output $Q_F$ whenever $\ex a>\ex f/4$, $\ex a+\ex f/2<1$, and
$\ex k_{\rm f}<\min(\ex a,\ex f,1-\ex a)$ (\cref{u:prop:fine-kl}).

This is the first place where a sharp estimate is needed. In
\cite[Proposition~5]{oai-typical-upper}, the relative entropy is at most
$2n^{1-l/2}$ for $0<l\le1/100$; the energy is computed on the event where the
empirical Gram matrix of the queries is close to the identity, and the
information is concentrated separately. Here the expected energy is computed
directly, through an exact conditional identity for the third moments of a
single query (\cref{u:lem:onequery}), and no Gram event enters. The saving
$\ex k_{\rm f}$ can then be taken arbitrarily close to $1/2$, and $1/2$ is the
limit of the estimate: the error $nh$ of the mesh and the number $n^{\ex a}$
of queries balance at $\ex a=1/2$. For $\ex a=\ex f=1/2$, we have
$\KL(Q_F\Vert\Pi_n)\le C_\beta n^{1/2+2\varepsilon_\omega}$, where
$\varepsilon_\omega=(1/4-\omega)/20$.

\paragraph{Coarse domination.} The second comparison is a bound on the
likelihood of the coarse proposal. There exists an event $E_B$ of the
coarse proposal, with $\PP(E_B^c)\le e^{-n^{1-\ex s_{\rm b}}}$, such that the
output restricted to $E_B$ has a kernel $q_B^{\rm good}(g,x)$ with respect to
$\gamma_n$ that satisfies
\begin{equation}\label{eq:ov-coarse-kernel}
 q_B^{\rm good}(g,x)\le e^{-nP_\beta+n^{1-\ex k_{\rm b}}}e^{\beta H^c_g(x)}
 =e^{n^{1-\ex k_{\rm b}}+F^c(g)-nP_\beta}\,\pi_g(x)
\end{equation}
for all $x$ and almost every $g$ (\cref{u:prop:coarse}). When
$F^c(g)\le nP_\beta+n^{1-\ex z_{\rm P}}$, this is a bound by
$e^{o(n^{1-\ex q})}$ times the Gibbs measure once $\ex q<\ex k_{\rm b}$ and
$\ex q<\ex z_{\rm P}$. Since the log density of $\Pi_n$ relative to $\gamma_n$
times the uniform law is $n\log2+\beta H^c_g(x)-F^c(g)$, it suffices to show
that, on $E_B$,
\begin{equation}\label{eq:ov-coarse-target}
 \frac1nI^{(n)}-\log2-\frac\beta nH^c_g(x)\le-P_\beta+n^{-\ex k_{\rm b}},
\end{equation}
where $I^{(n)}=\sum_iI_{V,i}$ is the logarithm of the density of the proposal
relative to the reference law of the paragraph on the forward entropy, at
$V=V_{\ex b}$. Integrating \eqref{eq:ov-coarse-target} against the reference
law, we obtain \eqref{eq:ov-coarse-kernel}; the factor $2^n$ cancels the
uniform spin weight $2^{-n}$.

The obstacle is the accuracy. By \eqref{eq:ov-scalar-identity}, the left
side of \eqref{eq:ov-coarse-target} has mean close to $-P_\beta$, and we must
show that it deviates by at most $n^{-\ex z}$, for an exponent
$\ex z>\ex k_{\rm b}>\ex q>\omega$, outside an event of probability
$e^{-n^{1-\ex s_{\rm b}}}$, with $\ex s_{\rm b}<1/4$. If the information and
the energy are concentrated separately, as in
\cite[Proposition~5]{oai-typical-upper}, each row has variance of order one,
and a deviation $n^{-\ex z}$ has probability of order
$e^{-n^{1-2\ex z}}$. In the parameter system below, the resulting requirement
$2\ex z<\ex s_{\rm b}$, together with $\omega<\ex q<\ex k_{\rm b}<\ex z$ and $\ex s_{\rm b}<1/4$ would
confine $\omega$ below $1/8$. We avoid this loss with three observations.

The first is algebraic. We write $H^c_g(x)/n=E_{\rm qry}+Z$, where $E_{\rm qry}$
is the energy carried by the queried columns and $Z$ collects the diagonal
entries $\langle u_j,Wu_j\rangle/n$ and the unread block. The queried energy is
an explicit function of the empirical Gram and cross matrices of the vectors
$\mathcal T_j$, $e_j$, and $x$, whose population values are known. Its
linearization in the errors of these matrices is itself an average of
independent rows: it equals $n^{-1}\sum_i\zeta_{h,i}$, where
\[
 \zeta_h=\sum_jw_j(A_{t_j}-A_{t_{j-1}})+\sum_js_j(x-m_{t_j})e_j-\frac{e_h}2,
 \qquad A_v=xm_v-\frac{m_v^2}2,
\]
for a row with spin $x$, magnetization $m$, and query endpoints $e_j$. The
deterministic weights are
$b_j=h_j^{-1/2}\int_{t_j}^{t_{j+1}}\rho(v)p(v)\dd v$ and $w_j=b_j/s_j$,
and $e_h=\sum_js_jb_j$ is the discrete scalar energy. With these
normalizations, $\E\zeta_h=e_h$ (\cref{u:lem:telescope}). The remainder of the linearization is quadratic in a
weighted measure $W_{\rm wt}$ of the errors (\cref{u:lem:weighted-energy}).
This gives
\[
 \frac1nI^{(n)}-\frac\beta nH^c_g(x)
 =\frac1n\sum_{i=1}^n\bigl(I_{V,i}-\beta\zeta_{h,i}\bigr)
 +O\bigl(D_*W_{\rm wt}^2\bigr)-\beta Z,
\]
where $D_*=1+\lceil\log_2N\rceil$.

The second observation is a pathwise cancellation between information and
energy (\cref{u:lem:joint}). On a mesh of vanishing size, $\zeta_h$ becomes
$\zeta_\infty=Z_V-e(V)/2$, where
$Z_V=\beta\int_0^Vp\dd A+\int_0^V\sqrt r\,(x-m)\dd B$. For every $V\ge v_*$,
we have exactly
\begin{equation}\label{eq:ov-joint}
 I_V-\beta\zeta_\infty=\log2-P_\beta-\frac{\beta^2}4p(V)^2
 +\frac{\beta^2p(V)}2(x-m_V)^2 .
\end{equation}
In other words, the stochastic integrals of the information and of the
energy cancel along every path. Note that the expectation of
\eqref{eq:ov-joint} is \eqref{eq:ov-scalar-identity}. To prove \eqref{eq:ov-joint}, we use It\^o's
formula for the Parisi solution along the scalar process
(\cref{u:lem:ito}). Before $v_*$, it shows that
$\beta^2pA+xY-U(v,Y)$ has differential $\beta\dd Z_v$, where $U$ is the
potential, and that the driver $B$ equals $M$; at $v_*$, the identity
$\log(1+x\tanh y)=xy-\log\cosh y$ matches the two sides; after $v_*$, the
difference $I_v-\beta Z_v$ moves only through $\beta^2p(v)(x-m_v)^2/2$. The
last term of \eqref{eq:ov-joint} does not have variance of order $h$: its
centered version $T_V$ has variance
$\beta^4p(V)^2\{4p(V)-3r(V)-p(V)^2\}/4\le\beta^4p(V)^3$, which is positive at
every fixed $V$. Its amplitude is at most $2\beta^2p(V)=O_\beta(n^{-\ex b/2})$,
however, and it is absorbed deterministically when $\ex z<\ex b/2$.

The third observation concerns the difference between the mesh and its limit.
The residual $R_h=Z_V-\zeta_h-e_h/2$ has mean $3(e(V)-e_h)/2=O_\beta(h)$, and
its centered moment generating function is at most $e^{C_\beta\theta^2h}$ for
$\abs\theta\le c_\beta/h$, uniformly in the horizon and in the number of cells
(\cref{u:lem:lag}). On the growing mesh, the lag of the magnetization inside a
cell of length $h_j$ contributes $h_j^2r(t_j)$, which is at most of order
$h^2$, and
there are at most $8/h$ cells. Since the rows are independent, we conclude that the average of
$I_{V,i}-\beta\zeta_{h,i}$ exceeds $\log2-P_\beta+C_\beta n^{-\ex z}$ with
probability at most $2e^{-c_\beta n^{1+\ex a-2\ex z}}$ (\cref{u:lem:row}). The
loss is $2\ex z-\ex a<0$, in place of $2\ex z$.

It remains to make $W_{\rm wt}^2$ smaller than $n^{-\ex z}$ and to control
$Z$. We have $W_{\rm wt}\le C_\beta n^{-\ex w}\log^Cn$ except with
probability
\[
 2\exp\bigl\{C_\beta n^{\ex a}-c_\beta n^{1-\max\{2\ex w,\ex w+\ex b/4\}}/\log^Cn\bigr\}
\]
(\cref{u:lem:weighted}). The Gram--Schmidt step is a Cholesky factorization,
and its errors involve the triangular truncation $\Upt$, which keeps the
strict upper triangle and half of the diagonal. As in
\cite[Section~4]{oai-typical-upper}, a dyadic decomposition of the triangle
costs only powers of $D_*$. The quantity $W_{\rm wt}$ measures the errors
through their projections onto dyadic intervals $I$ of query indices, weighted
by the population vectors. These projections are built from martingale
increments over the previous cells, and their variances are bounded by a
polylogarithmic multiple of the total length $d_I=\sum_{j\in I}h_{j-1}$ of
these cells. We divide each projection by $\sqrt{d_I}$ before applying a
variance-sensitive Bernstein inequality. Then the variance parameter is
polylogarithmic, and the inverse profile
$R_V\le C_\beta n^{\ex b/2}\sqrt{\log n}$ enters only the scale of the
exponential tail; this gives the loss $\ex w+\ex b/4$ beside $2\ex w$. The
half of the diagonal is a martingale plus a centered random drift, and both
have variance $O(d_I)$, with no power of $R_V$. At each dyadic depth, the
lengths $d_I$ add up to $t_N$, so that the squared errors, before the
normalization, add up without a
loss in the number of queries. The fixed Gram event
$\norm{\mathcal T^{\mathsf T}\mathcal T/n-I}\le(16D_*^2)^{-1}$, together with
its cross block, fails with probability whose exponent loses $\ex b/2$
(\cref{u:lem:gram}). Given the scalar paths and the spins, $Z$ is a centered
Gaussian variable, since the scalar recursion and the link never read the
diagonal entries or the unread block. Its conditional variance is at most
$C_\beta n^{-1}\{h+\tau^2+p(t_N)^2\}$ when $\norm{c-s}\le\tau$, where
$c_j=n^{-1}\langle u_j,x\rangle$ and $s_j$ is the population value of
$n^{-1}\langle\mathcal T_j,x\rangle$ (\cref{u:lem:weighted-energy}). The cruder bound $1/(2n)$ would
again lose $2\ex z$.

Collecting these estimates, we obtain \eqref{eq:ov-coarse-target} with
$\PP(E_B^c)\le e^{-n^{1-\ex s_{\rm b}}}$ whenever
\begin{gather*}
 \max\{2\ex w,\ex w+\ex b/4,\ex b/2,2\ex z-\ex a,2\ex z-2\ex w,2\ex z-\ex b\}
 <\ex s_{\rm b},\\
 \ex b/4<\ex a<1-\ex s_{\rm b},\qquad
 \ex k_{\rm b}<\ex z<\min(\ex a,\ex b/2,2\ex w).
\end{gather*}
This is the second place where a sharp estimate is needed. The failure bound
holds for the proposal law averaged over the disorder; no bound is asserted
for each fixed disorder.

\paragraph{Smoothing, coverage, and the connecting sequence.} A bound on
$\KL(Q_F\Vert\Pi_n)$ does not imply that the proposal reaches a given set of
appreciable Gibbs mass. To reverse the comparison, we follow
\cite[Section~5]{oai-typical-upper}. Conditionally on the spin $x$, the
disorder under $\Pi_n$ has density proportional to
$e^{-\norm g_F^2/2+\beta H^c_g(x)-F^c(g)}$, and since $F^c$ is convex and
$H^c_g(x)$ is linear in $g$, this law is $1$-strongly log-concave. The
transport inequality of Bogachev, Kolesnikov, and Medvedev
\cite[Corollary~3.10]{bogachev-kolesnikov-medvedev-2005} bounds the squared
Wasserstein distance from it by twice the relative entropy
(\cref{u:lem:transport}). We smooth the disorder by
$g\mapsto\sqrt{1-\delta}\,g+\sqrt\delta\,z$, with $\delta=n^{-\ex D}$ and an
independent standard Gaussian matrix $z$. As in
\cite[Lemma~11]{lugosi-neu-2022}, the transport bound turns the forward
entropy into a reverse entropy for the smoothed laws, at the cost of a factor
$\delta^{-1}$. Let $g_\star,g_\circ$ be independent with law $\gamma_n$, let
$g_{\rm base}=\sqrt{1-\delta}\,g_\star+\sqrt\delta\,g_\circ$, and let
$\nu^{\rm sm}_{g_{\rm base}}$ be the law of the spin of the fine proposal run
on the matrix $\sqrt{1-\delta}\,g_{\rm base}+\sqrt\delta\,z$, with fresh
auxiliary inputs. Then
\[
 \E\KL\bigl(\pi_{g_\star}\Vert\nu^{\rm sm}_{g_{\rm base}}\bigr)
 \le\frac{\beta^2(n-1)}2\delta+\delta^{-1}\KL(Q_F\Vert\Pi_n)
\]
(\cref{u:lem:smoothing}). Both terms are of order $n^{1-\ex D}$ and
$n^{1-\ex k_{\rm f}+\ex D}$, and the relative entropy is at most
$n^{1-\ex u_{\rm E}}$ with probability tending to one when
$\ex u_{\rm E}<\min(\ex D,\ex k_{\rm f}-\ex D)$. On this event, using the
inequality for the binary relative entropy, we find that, simultaneously for
all $S\subseteq\Sigma_n$,
\[
 \pi_{g_\star}(S)\ge n^{-\ex d}\quad\Longrightarrow\quad
 \nu^{\rm sm}_{g_{\rm base}}(S)\ge
 \exp\bigl\{-n^{\ex d}\bigl(n^{1-\ex u_{\rm E}}+\log2\bigr)\bigr\}
\]
(\cref{u:lem:coverage}). The set $S$ may depend on $g_\star$. Two independent
fine spins with law $\nu^{\rm sm}_{g_{\rm base}}$ fall on opposite sides of
every cut with both masses at least $n^{-\ex d}$ with probability
$\varpi_n\ge e^{-2n^{1-\ex e}}$ when $\ex e+\ex d<\ex u_{\rm E}$.

Suppose now that $g_\star\in\mathcal W_n$, and let $S_{g_\star}$ be the first
cut of small flow in a fixed order. The following argument is that of
\cite[Section~6]{oai-typical-upper}. If the two fine spins lie on opposite
sides of $S_{g_\star}$, we join them by a sequence of configurations with
short steps. Every such sequence comes close to the boundary of the cut. An
enlargement of the boundary has small Gibbs mass, and the likelihood bound
keeps the intermediate configurations out of it. We now describe the three
ingredients.

\emph{The excluded set.} Since $\norm{W(g_\star)}_{\rm op}\le12$, every
one-spin energy change is at most $24\sqrt n$, and every flip rate is at least
$e^{-24\beta\sqrt n}/2$. Then the vertex boundary of $S_{g_\star}$ has Gibbs mass
at most $4\phi_ne^{24\beta\sqrt n}$. We enlarge it by the Hamming radius
$R_n=\lceil n^{1-\ex r}\rceil$, and to compare Gibbs weights on the
enlargement, we also exclude configurations $x$ for which
\[
 \sum_i\bigl(\mathbf 1_A(i)-\mathbf 1_B(i)\bigr)
     (W(g_\star)x)_i\ge n^{1-\ex u_{\rm F}}
\]
for some $A,B\subseteq[n]$ with $|A|,|B|\le R_n$.
The expected Gibbs mass of the resulting set
$\mathcal E_n(g_\star)$ is at most
$e^{-n^{1-\ex q}/2}$ when $0<\ex q<\ex u_{\rm F}<\ex r<1$,
$2\ex u_{\rm F}-\ex r<\ex q$, and $\ex q<1/2$ (\cref{u:lem:boundary}). The
bound on the local fields uses the transport inequality for the conditional
disorder law given the spin. The cut is selected from $g_\star$ alone, and we
also restrict to the event $F^c(g_\star)\le nP_\beta+n^{1-\ex z_{\rm P}}$,
which fails with probability at most $e^{-c_\beta n^{1-2\ex z_{\rm P}}}$
(\cref{u:lem:pressure-tail}).

\emph{The sequence.} With independent standard Gaussian matrices $g_1,g_2$,
we set
\[
 g_\theta=\sqrt{1-\delta}\,g_{\rm base}
 +\sqrt\delta\,(\cos\theta\,g_1+\sin\theta\,g_2),\qquad
 0\le\theta\le\pi/2,
\]
and we rotate two independent copies of the other Gaussian inputs of the
proposal in the same way. For each fixed $\theta$, the matrix $g_\theta$ has
law $\gamma_n$, and the rotated inputs have their original law. We run the
coarse proposal at the $\ell_n+1$ points of a grid of spacing of order
$e^{-n^{\ex g_{\rm rot}}}$, and the fine proposal only at the two endpoints
$\theta=0$ and $\theta=\pi/2$. Given $(g_\star,g_\circ)$, the two fine spins
are independent with law $\nu^{\rm sm}_{g_{\rm base}}$. The sequence consists
of the first fine spin, the coarse spins in the order of the grid, and the
second fine spin. On an event whose complement has probability at most
$e^{-n^{1-\ex s_{\rm seq}}}$, the following holds (\cref{u:prop:sequence}):
\begin{itemize}
\item every coarse spin avoids $\mathcal E_n(g_\star)$: by the regression
$g_\theta=(1-\delta)g_\star+\sqrt{\delta(2-\delta)}\,Z_\theta$, with
$Z_\theta$ independent of $g_\star$, the conditional expectation of
$e^{\beta H^c_{g_\theta}(x)}$ given $g_\star$ is at most
$e^{O_\beta(n\delta)}e^{\beta H^c_{g_\star}(x)}$ when
$\norm{W(g_\star)}_{\rm op}\le12$, and
\eqref{eq:ov-coarse-kernel} bounds the probability of hitting the excluded
set by $e^{-n^{1-\ex q}/4}$ when $\ex q<\min(\ex k_{\rm b},\ex z_{\rm P},\ex D)$
(\cref{u:lem:transfer});
\item consecutive coarse magnetizations differ by at most
$e^{-n^{\ex g_{\rm rot}}/2}$ when $\ex a<\ex g_{\rm rot}$, because the
recursion amplifies differences of its inputs by at most
$e^{C_\beta n^{\ex a}\log n}$ on the Gram events at both ends of a grid step
(\cref{u:lem:stability});
\item every spin, coarse or fine, is within squared Euclidean distance $R_n/4$
of its coarse magnetization when $\ex r<\ex b/2$.
The fine and coarse runs share their paths up to $V_{\ex b}$;
the expected squared discrepancy per coordinate is
$p(V_{\ex b})\le C_\beta n^{-\ex b/2}$. The expected total
squared Euclidean distance is $np(V_{\ex b})=o(R_n)$;
independence across rows gives the tail in \cref{u:lem:pinning}.
\end{itemize}
On this event, we have that every step of the sequence has Hamming length at
most $R_n$, and every step contains a coarse spin. The union over the grid costs a factor
$e^{O(n^{\ex g_{\rm rot}})}$, which is absorbed when
$\ex g_{\rm rot}<1-\ex s_{\rm b}$, and the comparison of the failures requires
$\ex r<\ex s_{\rm b}$, $\ex q<\ex s_{\rm b}$, and
$\ex s_{\rm b}<\ex s_{\rm seq}$.

The fine proposal enters only through its relative entropy, which has no
exceptional event, and through the two endpoint spins, which are joined to the
grid by pinning. Every event in the union over the grid concerns the coarse
proposal, so that the sequence needs no Gram event of the fine proposal and
only the failure bound $e^{-n^{1-\ex s_{\rm b}}}$ of the coarse one.

\emph{The contradiction.} If the fine spins lie on opposite sides of
$S_{g_\star}$, some step of the sequence crosses the cut. A shortest path
across this step meets the boundary, so that the coarse endpoint of the step
lies within distance $R_n$ of the boundary, in $\mathcal E_n(g_\star)$. On
the event of the sequence, this is impossible. Given an outer pair
$(g_\star,g_\circ)$ in the event of the coverage bound, with $g_\star$ in
$\mathcal W_n$ and in the pressure event, we conclude that $\varpi_n$ is at most the
conditional probability that the sequence fails. Integrating over the outer
pair, and paying the coverage and pressure exceptions once, we find that
\begin{equation}\label{eq:ov-witness}
 \gamma_n(\mathcal W_n)\le
 C_\beta\bigl(n^{\ex u_{\rm E}-\ex D}+n^{\ex u_{\rm E}-\ex k_{\rm f}+\ex D}\bigr)
 +e^{-c_\beta n^{1-2\ex z_{\rm P}}}
 +e^{2n^{1-\ex e}-n^{1-\ex s_{\rm seq}}}
\end{equation}
(\cref{u:prop:witness}), which tends to zero when $\ex s_{\rm seq}<\ex e$.

\paragraph{The parameters.} The exponents of \cref{tab:ov-upper-exponents}
are affine functions of $\varepsilon_\omega=(1/4-\omega)/20\in(0,1/80)$, and
they satisfy every condition stated above, with strict inequalities, for
every $\omega\in(0,1/4)$ (\cref{u:prop:range}). The complete list of
conditions is given in \cref{u:sec:budget}. Four chains among them limit
the saving:
\begin{gather*}
 \ex k_{\rm f}<\min(\ex a,\ex f,1-\ex a)\le\frac12,\qquad
 \ex e+\ex d<\ex u_{\rm E}<\min(\ex D,\ex k_{\rm f}-\ex D)
 \le\frac{\ex k_{\rm f}}2,\\
 \ex q<\ex s_{\rm b}<\ex s_{\rm seq}<\ex e,\qquad
 \omega<\ex q .
\end{gather*}
They come from the forward entropy, from smoothing and coverage, from the
connecting sequence, and from \eqref{eq:ov-pruning}. Together they give
\[
 \omega<\ex q<\ex s_{\rm b}<\ex s_{\rm seq}<\ex e<\ex u_{\rm E}
 <\frac{\ex k_{\rm f}}2\le\frac14 .
\]
The conditions of coarse domination fit inside the third chain: the family
has $\ex q<\ex k_{\rm b}<\ex z<\ex b/2<\ex s_{\rm b}$, with consecutive gaps
$\varepsilon_\omega$, $\varepsilon_\omega$, $\varepsilon_\omega$, and
$2\varepsilon_\omega$, and a loss $2\ex z$ in place of $2\ex z-\ex a$ would
not fit. For $\omega=1/5$, we have $\varepsilon_\omega=1/400$,
the chain reads
\[
 \frac15<\frac9{40}<\frac{19}{80}<\frac6{25}<\frac{97}{400}<\frac{49}{200}
 <\frac{99}{400},
\]
and \eqref{eq:ov-witness} becomes
$C_\beta n^{-1/400}+e^{-c_\beta n^{2/5}}+\exp\{2n^{303/400}-n^{19/25}\}$.

Both arguments stop at the exponent $1/4$. The proof of
\cref{thm:barrier-730} needs $n\rho^4/\log n\to\infty$ for the overlap scale
$\rho$, and the proof of \cref{thm:upper} needs the chain of inequalities
$\omega<\ex q<\ex s_{\rm b}<\ex s_{\rm seq}<\ex e<\ex u_{\rm E}
<\ex k_{\rm f}/2\le1/4$ of \cref{u:sec:budget}. We do not know whether this
coincidence has a common explanation.


\section{Fluctuations and the mean bias}\label{sec:ov-fluct}

\Cref{f:main} is OpenAI's theorem \cite[Theorem~1.1]{oai-fluctuations}. We
give no new proof of it; our contribution concerning it and
\cref{cor:mean-bias} is described in \cref{sec:formal}. In
\cite{oai-fluctuations}, the theorem is reduced to a two-sided bound on a positive logarithmic moment. Let $F_n^c=F_n+\beta G$ be
the completed log partition function, and let $K_n(t)=\log\E\exp\{t(F_n^c-\E F_n^c)\}$.
For fixed $m\in(0,1/300)$ and $s_n=n^{-1/6+5m}$, it is proved in
\cite[Proposition~7.5]{oai-fluctuations} that
$cn^{20m}\le K_n(s_n)\le Cn^{100m}$ for all sufficiently large $n$
(\cref{f:window}); \cref{sec:fl-route} summarizes this proof with locators.
OpenAI bounds $K_n(s_n)$ by comparison with finite Gaussian hierarchies along
overlap quantiles. The error
of such a comparison is small when its minimizers satisfy a two-sided linear
bound. This bound follows from a one-site cavity estimate
\cite[Proposition~4.1]{oai-fluctuations}, from the asymptotic linearity of the
rescaled overlap quantiles of the minimizers
\cite[Theorem~5.1]{oai-fluctuations}, and from finite propagation through a bounded number of successive comparisons
\cite[Proposition~6.6]{oai-fluctuations}. The lower bound uses a chord
inequality at $\varepsilon=n^{-1/3}$, and the upper bound an iteration in
$\varepsilon$ with a low-field perturbation
\cite[Sections~7.1--7.4]{oai-fluctuations}.

A bound on $K_n$ at one tilt does not control $\Var(F_n)$ by itself, since a
rare deviation can dominate an exponential moment or a variance. The additional
input is that the sum of $F_n^c$ and a random variable with the exponential law
of rate one, independent of $(J,G)$, has a log-concave density
(\cref{f:logconcave}). This observation is
due to Aronow and Lopatto \cite[Proposition~5.1]{quantitative-parisi}, and
\cite[Lemma~7.6]{oai-fluctuations} states it in the present setting. For a
log-concave law with standard deviation $\sigma$, the centered logarithmic
moment is bounded at tilts below $c/\sigma$ and grows at least linearly in
$t\sigma$ at larger tilts $t$, and every interval of length $2r$ has
probability at most $Cr/\sigma$ (\cref{f:one-dim}). We deduce that
$\Var(F_n)=n^{1/3+o(1)}$ and, using the bound on intervals, the lower bound on
$\abs{F_n-\E F_n}$ in \eqref{f:typical} (\cref{sec:fl-logconcave,sec:fl-endgame},
following \cite[Sections~7.5 and~7.6]{oai-fluctuations}).

\Cref{cor:mean-bias} is the lower bound $n^{-5/6-\varepsilon}$ on the mean
bias. It follows from the variance--bias inequality
$\Delta_{n,\beta}\ge c_\beta((\Var(F_n)+\beta^2/2)/n)^{5/4}$ of Aronow and
Lopatto \cite[Lemma~7.4 and Remark~7.14]{quantitative-parisi}
(\cref{prop:bias-variance-deficit}). Inserting
$\Var(F_n)\ge n^{1/3-4\varepsilon/5}$ from \cref{f:main}, we find that
$\Delta_{n,\beta}\ge c_\beta n^{-5/6-\varepsilon}$ for all sufficiently large
$n$. Together with the upper bound $\Delta_{n,\beta}\le C_bn^{-7/10}$ of
\cref{inp:hb-rate}, which determines the exponent of \cref{thm:barrier-730}, this
shows that if $\Delta_{n,\beta}=n^{-a+o(1)}$ for some $a$, then
$7/10\le a\le5/6$. Both deductions are already recorded in
\cite[Section~1.6]{hexagon-bias}.
